% concatenated from main.tex, macros.tex
% SIAM Article Template
%\documentclass[review,onefignum,onetabnum]{siamonline250211}
\documentclass[onefignum,onetabnum]{siamonline250211}

% Information that is shared between the article and the supplement
% (title and author information, macros, packages, etc.) goes into
% ex_shared.tex. If there is no supplement, this file can be included
% directly.

% SIAM Shared Information Template
% This is information that is shared between the main document and any
% supplement. If no supplement is required, then this information can
% be included directly in the main document.


% Packages and macros go here
\usepackage{lipsum}
\usepackage{amsfonts}
\usepackage{graphicx}
\usepackage{epstopdf}
\usepackage{booktabs}
\usepackage{multirow}
\usepackage{subcaption}
\ifpdf
  \DeclareGraphicsExtensions{.eps,.pdf,.png,.jpg}
\else
  \DeclareGraphicsExtensions{.eps}
\fi

% Prevent itemized lists from running into the left margin inside theorems and proofs
\usepackage{enumitem}
\setlist[enumerate]{leftmargin=.5in}
\setlist[itemize]{leftmargin=.5in}

% Add a serial/Oxford comma by default.
\newcommand{\creflastconjunction}{, and~}

% Used for creating new theorem and remark environments
\newsiamremark{remark}{Remark}
\newsiamremark{hypothesis}{Hypothesis}
\crefname{hypothesis}{Hypothesis}{Hypotheses}
\newsiamthm{claim}{Claim}
\newsiamremark{fact}{Fact}
\crefname{fact}{Fact}{Facts}

% Sets running headers as well as PDF title and authors
\headers{Operator Learning for Smoothing and Forecasting}{E. Calvello, E. Carlson, N. Kovachki, M.N. Manta and A.M. Stuart}

% Title. If the supplement option is on, then "Supplementary Material"
% is automatically inserted before the title.
\title{Operator Learning for Smoothing and Forecasting}%\thanks{Submitted to the editors DATE.
%\funding{This work was funded by the Fog Research Institute under contract no.~FRI-454.}}}

% Authors: full names plus addresses.
\author{Edoardo Calvello\thanks{California Institute of Technology, Pasadena, CA 
  (\email{e.calvello@caltech.edu}).}
\and Elizabeth Carlson\thanks{Oregon State University, Corvallis, OR 
  (\email{carleliz@oregonstate.edu}).}
\and Nikola Kovachki\footnotemark[3]\thanks{NVIDIA Corporation, Santa Clara, CA 
  (\email{nkovachki@nvidia.com}).}
\and Michael N. Manta\thanks{Stanford University, Stanford, CA 
  (\email{mnmanta@stanford.edu}).}
\and \newline Andrew M. Stuart\thanks{California Institute of Technology, Pasadena, CA 
  (\email{astuart@caltech.edu}).}
  }

\usepackage{amsopn}
\DeclareMathOperator{\diag}{diag}

\usepackage{bbm}
%\usepackage[margin=1.0in]{geometry}
\usepackage{stmaryrd}
\usepackage[table,svgnames]{xcolor} % Required for \rowcolor
%\usepackage{booktabs}
\usepackage{hhline}
% Reals, complex, naturals
\newcommand{\1}{\mathbbm{1}}
\newcommand{\R}{\mathbbm{R}}
\newcommand{\Z}{\mathbbm{Z}}
\newcommand{\C}{\mathbbm{C}}
\newcommand{\N}{\mathbbm{N}}
\newcommand{\la}{\langle}
\newcommand{\ra}{\rangle}
\newcommand{\cH}{\mathcal{H}}
\newcommand{\cA}{\mathcal{A}}
\newcommand{\cN}{\mathcal{N}}
\newcommand{\cU}{\mathcal{U}}
\newcommand{\cV}{\mathcal{V}}
\newcommand{\cZ}{\mathcal{Z}}
\newcommand{\cG}{\mathcal{G}}
\newcommand{\cR}{\mathcal{R}}
\newcommand{\cX}{\mathcal{X}}
\newcommand{\cK}{\mathcal{K}}
\newcommand{\cE}{\mathcal{E}}
\newcommand{\cQ}{\mathcal{Q}}
\newcommand{\cD}{\mathcal{D}}
\newcommand{\cL}{\mathcal{L}}
\newcommand{\cO}{\mathcal{O}}
\newcommand{\cW}{\mathcal{W}}

\newcommand{\cT}{\mathfrak{T}}


\newcommand{\sA}{\mathsf{A}}
\newcommand{\sAp}{\mathsf{A}_{\mathsf{P}}}
\newcommand{\sC}{\mathsf{C}}
\newcommand{\sCp}{\mathsf{C}_{\mathsf{P}}}
\newcommand{\sQ}{\mathsf{Q}}
\newcommand{\sK}{\mathsf{K}}
\newcommand{\sL}{\mathsf{L}}
\newcommand{\sT}{\mathsf{T}}
\newcommand{\sF}{\mathsf{F}}
\newcommand{\sE}{\mathsf{E}}
\newcommand{\sS}{\mathsf{S}}
\newcommand{\sG}{\mathsf{G}}
\newcommand{\su}{\mathsf{u}}
\newcommand{\sv}{\mathsf{v}}
\newcommand{\rd}{{\mathrm d}}
\newcommand{\dK}{d_\textrm{\scalebox{.7}{$K$}}}


\newcommand{\sX}{\mathsf{X}}
\newcommand{\sY}{\mathsf{Y}}
\newcommand{\sV}{\mathsf{V}}
\newcommand{\sW}{\mathsf{W}}
\newcommand{\cF}{\mathcal{F}}
\newcommand{\cI}{\mathcal{I}}
\newcommand{\bbR}{\mathbb{R}}
\newcommand{\bbP}{\mathbb{P}}
\newcommand{\bbE}{\mathbb{E}}
\newcommand{\Cov}{\mathrm{Cov}}
\newcommand{\Var}{\mathrm{Var}}
\newcommand{\placeholder}{\mathord{\color{black!33}\bullet}}
\newcommand{\placeholderB}{\mathord{\color{black}\bullet}}
\newcommand{\range}[2]{\llbracket #1, #2 \rrbracket}
\renewcommand{\l}{\ell}
\newcommand{\lf}{\l_{f}}
\renewcommand{\lg}{\l_{g}}
\newcommand{\lh}{\l_{h}}
\newcommand{\nq}{\|_{\tiny{Q}}}
\newcommand{\np}{\|_{\tiny{P}}}


\newcommand{\dep}{\delta p}
\newcommand{\deq}{\delta q}
\newcommand{\ip}[2]{\left<#1,#2\right>}
\newcommand{\ipq}[2]{\left<#1,#2\right>_{\tiny{Q}}}
\newcommand{\ipp}[2]{\left<#1,#2\right>_{\tiny{P}}}
\newcommand{\lep}{\left(}
\newcommand{\rep}{\right)}


% \newcommand{\st}{\ \text{s.t.} \ }
\newcommand{\fa}{\ \forall \ }

\usepackage{amsmath}  % for \DeclareMathOperator macro
\DeclareMathOperator*{\argmax}{arg\,max}
\DeclareMathOperator*{\argmin}{arg\,min}

%=============
% Theorems, etc.
%=============

% Clever ref
% \usepackage{cleveref}

% Algorithm
\usepackage{algpseudocode}% http://ctan.org/pkg/algorithmicx

\usepackage{amsmath}
\usepackage{amssymb}
\usepackage{amsfonts}
\usepackage{dsfont}
\usepackage{mathtools}


%%%%%%%%%%% From CAMS PAPER %%%%%%%%%%%%%%
%     If you need symbols beyond the basic set, uncomment this command.
\usepackage{amssymb}


% Recommended, but optional, packages for figures and better typesetting:
\usepackage{microtype}

% \usepackage{subcaption}
% \usepackage{subfigure}
% see: https://tex.stackexchange.com/questions/548735/latex-proper-ref-to-sub-figures-and-proper-alignment

%\usepackage[width=1.0\textwidth]{caption}

%\usepackage{booktabs} % for professional tables

% hyperref makes hyperlinks in the resulting PDF.
% If your build breaks (sometimes temporarily if a hyperlink spans a page)
% please comment out the following usepackage line and replace
% \usepackage{icml2021} with \usepackage[nohyperref]{icml2021} above.
% \usepackage{hyperref}

% Attempt to make hyperref and algorithmic work together better:
%\newcommand{\theHalgorithm}{\arabic{algorithm}}

% Use the following line for the initial blind version submitted for review:
% \usepackage[accepted]{icml2021}

% Coloring For Changes
% \usepackage{color}
% \usepackage{xcolor}
\usepackage[table,svgnames]{xcolor}         % colorse1

\usepackage{soul}

\definecolor{darkred}{rgb}{.7,0,0}
\definecolor{grey}{rgb}{.7,.6,.5}

\definecolor{darkgreen}{rgb}{0.05, 0.5, 0.06}
\newcommand{\ecal}[1]{{\color{purple}{#1}}}

\definecolor{darkred}{rgb}{0.5, 0.05, 0.06}
\newcommand{\as}[1]{{\color{darkred}{#1}}}

\newcommand{\ecar}[1]{{\color{cyan}{#1}}}
\newcommand{\mm}[1]{{\color{blue}{#1}}}
\newcommand{\nk}[1]{{\color{darkgreen}{#1}}}
% \newcommand{\mlrm}[1]{\textcolor{blue}{\sout{#1}}}

% useful colors for late-stage editing
\newcommand{\done}[1]{{\color{grey}{#1}}}
\newcommand{\agree}[1]{{\color{black}{#1}}}
\newcommand{\gottado}[1]{{\color{green}{#1}}}

\usepackage[normalem]{ulem}
\newcommand\mlremove{\bgroup\markoverwith
{\textcolor{red}{\rule[.5ex]{2pt}{1pt}}}\ULon}
\newcommand\asremove{\bgroup\markoverwith
{\textcolor{cyan}{\rule[.5ex]{2pt}{1pt}}}\ULon}
\newcommand\agreeremove{\bgroup\markoverwith
{\textcolor{blue}{\rule[.5ex]{2pt}{1pt}}}\ULon}


% BEGIN MATTs STUFF
\usepackage{amsmath,amsbsy,amsgen,amscd,amsfonts} %amsthm
% \usepackage{cleveref}
% \usepackage{bbm} % this helps get a nice indicator function with \mathbbm{1}
\usepackage{mathrsfs} % using this for \mathscr{g}
\usepackage{bbold}
% \usepackage[utf8]{inputenc}
\usepackage{csquotes}

%==============
% Symbols
%==============
%==============
% Constants
%==============

% Set constants upright
\newcommand{\cnst}[1]{\mathrm{#1}}
\newcommand{\econst}{\mathrm{e}}

\newcommand{\zerovct}{\vct{0}} % Zero vector
\newcommand{\Id}{{I}} % Identity matrix
\newcommand{\onemtx}{\bm{1}}
\newcommand{\zeromtx}{\bm{0}}

\newcommand{\ASelf}{A_\textrm{self}}
\newcommand{\ACross}{A_\textrm{cross}}
\newcommand{\sm}{\sigma_\textrm{softmax}}

\newcommand{\enc}{\text{enc}}
\newcommand{\dec}{\text{dec}}
\newcommand{\init}{\text{init}}

%==============
% Sets
%==============
\providecommand{\mathbbm}{\mathbb} % In case we don't load bbm

\newcommand{\sR}{\mathsf{R}}
\newcommand{\sD}{\mathsf{D}}
\newcommand{\fp}{\mathfrak{p}}
\newcommand{\fq}{\mathfrak{q}}
\newcommand{\ft}{\mathfrak{t}}
\newcommand{\fx}{\mathfrak{x}}
\newcommand{\tp}{\widetilde{p}}
\newcommand{\tq}{\widetilde{q}}

\newtheorem{assumption}{Assumption}[section]
\newtheorem{example}{Example}[section]

\newlist{contributions}{enumerate}{5}
\setlist[contributions]{font={\bfseries}}
\crefname{contributionsi}{Contribution}{Contributions}

\newlist{probenum}{enumerate}{5}
\setlist[probenum]{font={\bfseries}}
\crefname{probenumi}{Problem}{Data Assimilation Problems}

%==============
% Probability
%==============
\newcommand{\Prob}{\operatorname{\mathbbm{P}}}
\newcommand{\Expect}{\operatorname{\mathbb{E}}}

%==============
% Vectors and matrices
%==============
\newcommand{\vct}[1]{{#1}}
\newcommand{\mtx}[1]{{#1}}
\newcommand{\vd}{v^{\dagger}}
\newcommand{\pd}{p^{\dagger}}
\newcommand{\qd}{q^{\dagger}}
\newcommand{\Psid}{\Psi^{\dagger}}
\newcommand{\Gd}{G^{\dagger}}


%%% Local Variables: 
%%% mode:latex
%%% TeX-master: "ex_article"
%%% End: 


% Optional PDF information
\ifpdf
\hypersetup{
  pdftitle={An Example Article},
  pdfauthor={D. Doe, P. T. Frank, and J. E. Smith}
}
\fi

% The next statement enables references to information in the
% supplement. See the xr-hyperref package for details.

%\externaldocument[][nocite]{ex_supplement}

% FundRef data to be entered by SIAM
%<funding-group specific-use="FundRef">
%<award-group>
%<funding-source>
%<named-content content-type="funder-name"> 
%</named-content> 
%<named-content content-type="funder-identifier"> 
%</named-content>
%</funding-source>
%<award-id> </award-id>
%</award-group>
%</funding-group>

\begin{document}

\maketitle
% REQUIRED
\begin{abstract}
Machine learning has opened new frontiers in purely data-driven algorithms for data assimilation in,
and for forecasting of, dynamical systems; the resulting methods are showing some promise.
However, in contrast to model-driven
algorithms, analysis of these data-driven methods is poorly developed. 
In this paper we address this issue, developing a theory to underpin
data-driven methods to solve smoothing problems arising in data assimilation
and forecasting problems. The theoretical framework relies on two key components: (i) establishing
the existence of the mapping to be learned; (ii) the properties of 
the operator learning architecture used to approximate this mapping. 
By studying these two components in conjunction, we establish novel
universal approximation theorems for purely data-driven algorithms for both smoothing
and forecasting of dynamical systems. We work in the continuous time setting, hence deploying
neural operator architectures. The theoretical results are illustrated
with experiments studying the Lorenz `63, Lorenz `96 and Kuramoto-Sivashinsky
dynamical systems.
\end{abstract}

% REQUIRED
\begin{keywords} Data Assimilation, Smoothing, Forecasting, 
Neural Operators, Universal Approximation.  
\end{keywords}

% REQUIRED
%\begin{MSCcodes}

%\end{MSCcodes}

\section{Introduction}
\label{sec:intro}



Machine learning has opened new frontiers in data assimilation and forecasting in dynamical systems,
ranging from methods which improve on traditional model-driven 
algorithms, through to new purely data-driven methods. The goal
of this paper is to establish a mathematical approach to the study
of the purely data-driven methods. To this end, we consider the general dynamical system given by
\begin{subequations}
\label{eq:dynamical_system_intro}
\begin{align}
        \label{eq:dynsys_f}
        \dot{p} &= f(p, q), \qquad p(0) = p_0 \in \R^{d_p}, \\
        \label{eq:dynsys_g}
        \dot{q} &= g(p, q), \qquad q(0) = q_0 \in \R^{d_q}.
\end{align}
\end{subequations}
Our focus is on two problems: (i) existence, universal approximation, and approximation from data
of the map from observed $\{p(t)\}_{t\in [0,T]}$ to unobserved $\{q(t)\}_{t\in [0,T]}$; (ii)
existence, universal approximation, and approximation from data of the map from  observed $\{p(t)\}_{t\in [0,T]}$ to unobserved $\{p(t)\}_{t\in [T,T+\tau]}$.

\subsection{Context and Literature Review}
\label{sec:CLR}

In the context of \eqref{eq:dynamical_system_intro} there are three problem classes which naturally arise: (I) \textit{smoothing}  concerns the offline denoising and estimation of the observed and unobserved  states over some time interval, given observational data in that same time 
interval; an instance of this, which we study in this paper, is the estimation of unobserved 
$\{q(t)\}_{t\in [0,T]}$ from observed $\{p(t)\}_{t\in [0,T]}$; (II) \textit{forecasting} concerns estimation of a state over some future time interval, given observational data over a past time interval; an instance of this, which we also study in this paper, is the estimation of 
$\{p(t)\}_{t\in [T,T+\tau]}$ from observed $\{p(t)\}_{t\in [0,T]}$. Also of importance
is (III) \textit{filtering}, which concerns the online estimation of an underlying state, 
sequentially as observations are received; this, however, is not a focus of the current paper.

The field of data assimilation (DA) is concerned with the use of 
observational data to perform state estimation in dynamical systems (including the estimation of quantities that are 
not observed). Both smoothing and filtering are used in practice, often in conjunction
with forecasting, as in the field of weather prediction. Traditional model-driven techniques combine the observations with knowledge of the underlying dynamical
system to address the three tasks (I--III) \cite{kalnay2002atmoshperic,asch2016data,Evensenetal2022}.
Although primarily developed for geophysical applications in the ocean-atmosphere sciences,
such model-driven methods to address tasks (I--III) have been systematized and now constitute general-purpose methodologies
\cite{reich2015probabilistic,pandya2022review, sanz2023inverse}. In the last few years new purely data-driven algorithms, which 
do not require knowledge of the dynamical system when deployed, have started to emerge \cite{kurth2023fourcastnet,bi2023accurate,price2025probabilistic,Allen2025,alexe2024graphdop,kossaifi2026demystifying}. These methods have also been developed primarily in weather
forecasting, but present opportunities for deployment in many other domains. Developing
mathematical theory pertinent to these methods can play an important role
in the process of making the methods more widely applicable and is the focus of this paper.
 

Over the last half century, model-driven filtering, smoothing and forecasting methods have dominated 
in most application domains, starting from the seminal work of Kalman and Bucy for linear Gaussian
systems \cite{kalman1960new,kalman1961new}; the field evolved with introduction of the bootstrap, extended (ExKF) and ensemble Kalman filters (EnKF) \cite{doucet2000sequential,jazwinski2007stochastic,evensen1994sequential} which, respectively,
typically work well for small-, medium- and large-scale dynamical systems: bootstrap particle filter weights collapse in high dimensions and so they are best in low dimensions; EnKF has equal weights and has performed well for weather forecasting with state spaces in the billions of variables; ExKF falls between these cases as it makes a Gaussian approximation (mitigating weight collapse), but the covariances cannot be propagated in high dimension as they are too large. All model-driven approaches require knowledge of the dynamics, requiring
computationally expensive evaluations in the context of many large-scale applications of interest. Purely data-driven, model-free, methods for forecasting were proposed by Lorenz in 1969 \cite{lorenz1969atmospheric}, and go by the name of analog forecasting;
these methods were not widely adopted, however, primarily because of the lack 
of data to support them, and because they are discontinuous as a function of input data.
In the last decade kernel analog forecasting
has been developed \cite{alexander2020operator}, a smoothening of the original
approach of Lorenz, supported by theory. Furthermore, methods such as dynamic
mode decomposition \cite{schmid2010dmd}, which are also purely data-driven, 
have been widely adopted and come with a developing theory \cite{mezic2020spectrum}. These methods have not yet, however, found success in the context of large-scale data assimilation and forecasting problems; their success
is mainly in the identification of large-scale coherent features in high-dimensional systems.

Machine learning has recently emerged as a novel computational tool for 
improving data assimilation and forecasting in both the model- and data-driven settings.
Learning has been introduced within model-driven DA in a variety of ways,
primarily in the filtering context. For example, it has been used for model error correction \cite{Farchi2021,levine2022}, for constructing cheap surrogates of the dynamics \cite{sanzalonso2025longtime, adrian2025data} and more recently to approximate conditioning in the assimilation step \cite{Bach2025Learning}. A data-driven approach to smoothing has recently been employed in \cite{Allen2025, gupta2026healda} for estimating the initial conditions of weather dynamics. Machine learning has also introduced a new approach to forecasting which is fully data-driven; this involves the learning of mappings which output predictions for the forecasted components of the system from their past observations. This approach has garnered wide interest in domains like weather forecasting, where both model evaluations and assimilation steps are particularly costly \cite{kurth2023fourcastnet,bi2023accurate, price2025probabilistic, alexe2024graphdop, kossaifi2026demystifying}. Time-series forecasting for
commercial and economic purposes has also seen considerable impact from machine-learning
\cite{wen2017multi,eisenach2020mqtransformer,challu2023nhits,liu2023itransformer,ansari2024chronos}. In a similar context, deploying pretrained time-series foundation models for forecasting of partially observed chaotic dynamical systems has been numerically investigated in \cite{zhang2025zeroshot}. Underpinning these data-driven approaches
with theoretical foundations is the goal of this paper.
In \Cref{tab:da_comparison} we summarize the distinction between model-driven DA, that has dominated
for the last half century, and recently emerging data-driven approaches including, but not limited to, the specific methods we study in this paper.

\begin{table}[htbp]
    \centering
    \footnotesize
    \caption{Comparison between model-driven and data-driven approaches.}
    \label{tab:da_comparison}
    \begin{tabular}{lcc}
        \toprule
        & \textbf{Model-driven} & \textbf{Data-driven} \\
        \midrule
        \textbf{Algorithm} 
        & Uses dynamical physics model
        & Uses data from dynamical model \\ & & and/or direct observations \vspace{1ex}\\
        \textbf{Limitation} 
        & Possibly expensive to evaluate and
        & Requires existence of direct mapping, \\
        & may miss some dynamics due to & a sufficiently dense training dataset and may not  \\
        & unmodeled physics & enforce explicit physical constraints \vspace{1ex}\\
        \textbf{Advantage} 
        & Interpretable and output satisfies 
        & Possibly cheap to evaluate and \\
        & prescribed physical constraints & agnostic to model dynamics \\ 
        \bottomrule
    \end{tabular}
\end{table}

\iffalse
\begin{table}[htbp]
    \centering
    \footnotesize
    \caption{Comparison between model-driven and data-driven approaches.}
    \label{tab:da_comparison}
    \begin{tabular}{lcc}
        \toprule
        & \textbf{Model-driven} & \textbf{Data-driven} \\
        \midrule
        \textbf{Algorithm} 
        & Uses dynamical model
        & Does not use dynamical model \\
        \textbf{Limitation} 
        & Possibly expensive model evaluations
        & Requires existence of direct mapping \\
        \textbf{Advantage} 
        & Interpretable 
        & Agnostic to dynamics \\
        \bottomrule
    \end{tabular}
\end{table}
\fi 

\subsection{Contributions and Outline}
\label{subsec:contributions}

In this paper we develop theory for, and insights into,  
the data-driven smoothing and forecasting problems previously defined. We work in the
noise-free setting, and in continuous time,
focusing on three essential ingredients: (a) the interplay 
between the theoretical properties of the dynamics and the existence of well-defined maps
that can in principle be approximated; (b) the universal approximation 
properties of neural operators \cite{kovachki2023neural}, neural network architectures parametrizing mappings between spaces of functions, for these maps; and (c) the implementation and properties of approximations of
these maps learned from data in a model-free setting. We show that it is possible to construct the
desired operators, defined locally, under an observability-rank condition on the dynamics. We show that this condition is intimately connected to the setup of \cite{hermann1977nonlinear}, a foundational work in control theory which 
establishes that, under a more general observability-rank condition, 
distinct observations correspond to distinct dynamics. Takens' delay embedding theorem \cite{takens2006detecting,sauer1991embedology} is also conceptually pertinent to our work, but the concrete control-theoretic perspective of \cite{hermann1977nonlinear} provides a more natural basis for the theory and
algorithms that we develop here.
The existence of the desired operators allows approximation via neural operators, for which an emerging universal approximation theory has been developed \cite{lanthaler2023nonlocal}. In particular, we show how architectures such as the transformer neural operator \cite{calvello2024accuracy} may be applied in the smoothing and forecasting contexts, and 
yield desirable universal approximation results. 
We then implement these neural operators, demonstrating
data-driven approximation of the maps in practice.
Our contributions are as follows:
\vspace{1ex}
\begin{contributions}[label=(C\arabic*)]
\item 
\label{contrib:control}
We bridge the control theoretic literature with data-driven smoothing and forecasting by using an observability condition on the dynamics, which guarantees existence of a locally-defined continuous operator mapping an observed component of the dynamical system to an unobserved component. This condition is defined in \Cref{assump:observability}.
\vspace{1ex}
\item \label{contrib:approx_smoothing} We prove a theorem establishing the existence of a class of neural operators capable of approximating, to arbitrary accuracy, the operator mapping an observed component of the system to an unobserved component. This is the first universal approximation theorem for a smoothing problem in DA. The result is stated in \Cref{thm:UAsmoothing}.
\vspace{1ex}
\item \label{contrib:approx_forecasting} We prove a theorem establishing the existence of a class of neural operators capable of approximating, to arbitrary accuracy, the operator mapping a trajectory of the observed component of the system to the future of that trajectory. The result is stated in \Cref{thm:UAforecasting}.
\vspace{1ex}
\item \label{contrib:numerics}
We showcase the universal approximation theory developed by applying transformer neural operators on smoothing and forecasting problems in the setting of the Lorenz `63, Lorenz `96, and Kuramoto-Sivashinsky dynamical systems.
\end{contributions}
\vspace{1ex}
To the best of our knowledge, the Theorems \ref{thm:UAsmoothing} and \ref{thm:UAforecasting} are the first of their kind. After introducing notation that will be used throughout the paper in \Cref{subsec:notation}, we describe the general dynamical system setting in \Cref{subsec:setup}; here we also define the smoothing and forecasting problems which will be the object of investigation. In \Cref{sec:Observability} we introduce the observability framework
that underpins our analysis, leading to \Cref{contrib:control}. Building on regularity and observability assumptions on the dynamical systems, we develop the universal approximation theory for smoothing and forecasting in \Cref{sec:UAforDA}, thus addressing \Cref{contrib:approx_smoothing,contrib:approx_forecasting}. In \Cref{sec:experiments} we demonstrate this theory numerically, thus addressing \Cref{contrib:numerics}. Finally in \Cref{sec:conclusions} we discuss the results of the paper and conclude by outlining avenues for further work. 

\subsection{Notation}
\label{subsec:notation}

Throughout we denote the positive integers and non-negative integers
respectively by $\N=\{1,2,\cdots\}$ and $\Z^+=\{0,1,2,\cdots\},$
and the notation $\R=(-\infty,\infty)$ and $\R^+=[0,\infty)$ for the
reals and the non-negative reals. For a set $D \subset \R^m$, we denote 
by $\Bar{D}$ the closure. 
Given two vector spaces $\cU,\cV$ we denote by $\sL(\cU;\cV)$ the space of linear operators acting between $\cU$ and $\cV$. We denote by $C(D;\R^d)$ the infinite dimensional Banach space of continuous functions mapping the set $D$ to the $d$-dimensional vector space $\R^d$. The space is endowed with the supremum norm $\|\placeholder\|_\infty$. We will sometimes use the shorthand notation $C(D)$ to denote continuous functions defined on the domain $D$,
whenever the image space is unambiguous or not germaine. Similarly, we sometimes drop the domain notation completely, when it is clear from context. For $s \in \Z^+$ we denote by $C^s$ the space of continuously differentiable functions up to order $s$. Let $C^{\infty}$ denote the infinite dimensional real vector space of all infinitely differentiable functions. For any function $f\in C^s\bigl(\R^{d_1};\R^{d_2}\bigr)$ and any $0\leq k\leq s$, we define by $\sD^k f(v)$ the $k$th derivative at $v\in \R^{d_1}$ viewed as the linear operator $\sL\bigl(\bigotimes_{i=1}^k\R^{d_1}; \R^{d_2}\bigr)$. We endow the space $C^s$ with the norm $\|\placeholder\|_{C^s}$ defined as
\begin{equation}
\label{eq:norm}
    \|f\|_{C^s} = \sum_{j=0}^s \| \sD^j f\|_\infty,
\end{equation}
for any function $f\in C^s$. Given $s \in \N$, for any function $f\in C^s\bigl(\R^{d_1};\R^{d_2}\bigr)$ and a vector field $w\in C^{s-1}\bigl(\R^{d_1};\R^{d_1}\bigr)$, we also define $\mathcal{L}_w\,f(v)$, the Lie derivative of $f$ along $w$, as
\[
\mathcal{L}_w\,f(v) = \sD f (v) w(v), \qquad v\in \R^{d_1}.
\]
Note that $\mathcal{L}_w\,f \in C^{s-1}\bigl(\R^{d_1};\R^{d_2}\bigr).$
For $k\leq s$ we also denote by $\cL^k_w \,f$ the Lie derivative of $f$ along $w$ of order $k$, defined by $k-$fold application of $\cL_w$. We will occasionally drop the vector field from the notation when it is clear 
from the context.

\subsection{Dynamical System Setup}
\label{subsec:setup}



Consider the dynamical system \eqref{eq:dynamical_system_intro}.
We make the following regularity assumption:




\begin{assumption}[Regularity]
\label{assump:regularity}
 Assume that, for some $k \in \N$,
 \begin{itemize}
     \item $f \in C^k (\R^{d_p + d_q}; \R^{d_p})$;
     \item $g \in C^k (\R^{d_p + d_q}; \R^{d_q})$.
 \end{itemize}
Assume, furthermore, that solutions to \eqref{eq:dynamical_system_intro} exist for all $t \in \R^+.$
\end{assumption}

Under \Cref{assump:regularity} it follows that, for any $T>0$, $p \in C^k ([0,T];\R^{d_p})$ and $q \in C^k ([0,T];\R^{d_q})$. We define by $\Phi:\R^+ \times\R^{d_p+d_q}\to\R^{d_p+d_q}$ the semigroup of operators associated to the dynamical system, so that $\bigl(p(t),q(t)\bigr)=\Phi(t,p_0,q_0).$  
Let $I \subset \R^{d_p + d_q}$ denote a compact set of initial conditions to \eqref{eq:dynamical_system_intro} and fix $T>0$. Since by \Cref{assump:regularity} solutions to \eqref{eq:dynamical_system_intro} exist for all $t \in \R^+$, the result of \cite[Theorem 2.10]{teschl2012ordinary} implies that $\Phi(\placeholder,\placeholder)\in C^k([0,T]\times I;\R^{d_p + d_q})$; therefore, it may be deduced that the map $x\mapsto \Phi(\placeholder,x)\in C^k([0,T];\R^{d_p + d_q})$ is continuous. Since the image of a compact set under a continuous map is compact, we deduce that the set of orbits 
$S^I_{[0,T]} = \{\Phi(t,x)\,\text{for }t\in [0,T]: x \in I \}$ is compact in
$C^k([0,T];\R^{d_p + d_q})$. Let $\pi_p \colon \R^{d_p + d_q} \to \R^{d_p}$ be the projection map onto the first $d_p$ coordinates and $\pi_q \colon \R^{d_p + d_q} \to \R^{d_q}$ be the projection map onto the remaining $d_q$ coordinates of the state space. We note that these are both continuous functions. Define the projected sets
\begin{equation}
\label{eq:Ssets}
S^I_{[0,T],p} = \{\pi_p (s): s \in S^I_{[0,T]}\}, \quad S^I_{[0,T],q} = \{\pi_q (s): s \in S^I_{[0,T]}\};
\end{equation}
these too are compact in $C^k([0,T];\R^{d_p})$ and $C^k([0,T];\R^{d_q})$ since $S^I_{[0,T]}$ is compact and projection is continuous. Fixing some $\tau>0$, for the set of orbits defined on the future time interval $S^I_{[T,T+\tau]} = \{\Phi(t,x)\,\text{for }t\in [T,T+\tau]: x \in I \}$ we similarly define the projected set
\begin{equation}
    \label{eq:forecastingSet}
    S^I_{[T,T+\tau],p} = \{\pi_p (s): s \in S^I_{[T,T+\tau]}\}.
\end{equation}
It is readily deduced that $S^I_{[T,T+\tau],p}$ is also compact in $C^k([T, T+\tau];\R^{d_p})$. 
Given the dynamical systems \eqref{eq:dynamical_system_intro} equipped with the regularity assumption in \Cref{assump:regularity}, we define the two problems central to this paper. We assume that
pair $(p,q)$ solves \eqref{eq:dynamical_system_intro} but that vector field $(f,g)$ is not known to us. We then consider the following problems.
\vspace{1ex}
\begin{probenum}[label=(P\arabic*)]

\item \label{prob:smoothing} \textit{Smoothing}: recover $q\in S^I_{[0,T],q}\subset C^k ([0,T];\R^{d_q})$ from observed $p\in S^I_{[0,T],p} \subset C^k ([0,T];\R^{d_p})$. % for $(p, q) \in S^I_{[0,T]}$.
\vspace{1ex}
\item \label{prob:forecasting} \textit{Forecasting}: predict
$p \in S^I_{[T,T+\tau],p}$ from $p \in S^I_{[0,T],p}.$ 


\end{probenum}



\section{Observability}
\label{sec:Observability}
%\label{sec:UAforDA}


In this section we develop the observability framework that underpins the existence and universal
approximation of smoothing and forecasting maps, the subject of \Cref{sec:UAforDA}. 
Consider dynamical system \eqref{eq:dynamical_system_intro}. Underlying the theory in \Cref{sec:UAforDA}
is the question of whether a map $\{p(t)\}_{t\in [0,T]} \mapsto \{q(t)\}_{t\in [0,T]}$ exists. If this
map exists then the maps defined by the smoothing and forecasting problems both exist. A sufficient condition for the map $\{p(t)\}_{t\in [0,T]} \mapsto \{q(t)\}_{t\in [0,T]}$ to exist 
is that initial condition $q_0$ is determined by $\{p(t)\}_{t\in [0,T]}.$ Then, equation \eqref{eq:dynsys_g} can be integrated as a non-autonomous equation for $q(\cdot)$, driven by the observed $p(\cdot)$, defining the smoothing map. Once this is well-defined, $(p(T),q(T))$ are known and integration of \eqref{eq:dynamical_system_intro} defines the forecasting map.


This section is hence focused on the question of 
recovering $q_0$ from $p$ and its derivatives at time $t=0,$ noting that these derivatives are determined by
$\{p(t)\}_{t\in [0,T]}$, assuming sufficient differentiability. In fact we formulate the question a bit
more generally, seeking to recover $q(t)$ at some given time $t$, given $p$ and its derivatives at that time.
In \Cref{ssec:noto} we introduce the notation required to discuss the observability rank condition
which encapsulates solvability for $q(t)$ at some time $t$.  The observability-rank condition itself is
introduced in \Cref{ssec:obs}. We conclude in \Cref{ex:Lorenz} with an example of
observability in the Lorenz `63 equation.

\subsection{Observability Setup}
\label{ssec:noto}

The results in \Cref{sec:UAforDA} are obtained via an underlying observability assumption on the system \eqref{eq:dynamical_system_intro}. To state such an observability condition we introduce the following notation relating to the dynamical system. Let $\cL^m$ denote the $m^{th}$ Lie derivative of $f$ along the vector field $(f,g)$ and let $\widetilde{F}^{(n)}\colon\R^{d_p+d_q}\to\R^{nd_p}$ be defined from the first $n$ Lie derivatives of $f$ along the vector field $(f,g)$ as follows:
    \begin{equation}
    \label{eq:Lie_tilde}
    \widetilde{F}^{(n)}(\mathfrak{p},\mathfrak{q}) = \Bigl( \cL^0\,f(\mathfrak{p},\mathfrak{q})^\top,\cL^1\,f(\mathfrak{p},\mathfrak{q})^\top,\ldots, \cL^{n-1}\,f(\mathfrak{p},\mathfrak{q})^\top\Bigr)^\top,
    \end{equation}
    for any $(\fp, \fq) \in \R^{d_p + d_q}$.
We also define $F^{(n)}\colon\R^{d_p+d_q}\to\R^{(n+1)d_p}$ by
    \begin{equation}
    \label{eq:Lie}
    F^{(n)}(\mathfrak{p},\mathfrak{q}) = \Bigl(\fp^{\top},\, \cL^0\,f(\mathfrak{p},\mathfrak{q})^\top,\cL^1\,f(\mathfrak{p},\mathfrak{q})^\top,\ldots, \cL^{n-1}\,f(\mathfrak{p},\mathfrak{q})^\top\Bigr)^\top,
    \end{equation}
    for any $(\fp, \fq) \in \R^{d_p + d_q}$. We note that the two functions $\widetilde{F}^{(n)}$ and
    $F^{(n)}$ differ only through the inclusion of $\fp\in \R^{d_p}$ in $F^{(n)}.$ Using either we can write down equations which, potentially, enable the extraction of $q(t)$ from knowledge of $p(\cdot)$ and its derivatives at $t.$  
    
    
    To this end we show how a collection of Lie derivatives of the vector field $f$ can be expressed in terms of derivatives of $p.$ By considering \eqref{eq:dynsys_f}, we may deduce that for any $(p,q) \in S^I_{[0,T]},$ it holds that
    \begin{align*}
    \cL\,f \bigl(p(t),q(t) \bigr) &= \partial_p f\bigl(p(t),q(t) \bigr)\cdot f \bigl(p(t),q(t) \bigr) + \partial_q f\bigl(p(t),q(t) \bigr)\cdot g \bigl(p(t),q(t) \bigr)\\
    & = \partial_t^2\, p(t),
    \end{align*}
    for all $t\in[0,T]$.
    By repeated application of Lie derivatives, it holds for any $(p, q)\in S^I_{[0,T]} \subset C^k([0,T];\R^{d_p})\times C^k([0,T];\R^{d_q})$ and for all $n \leq k-1$ that
    \begin{equation}
    \label{eq:connection_control}
    \cL^{n}\,f \bigl(p(t),q(t) \bigr) = \partial_t^{n+1}\, p(t),
    \end{equation}
    for all $t\in [0,T]$. Define operator $\widetilde{P}^{(n)}\colon C^k([0,T];\R^{d_p})\to \bigotimes_{j=1}^{n} C^{k-j}([0,T];\R^{d_p})$ via its action on $p\in C^k([0,T];\R^{d_p})$ as follows:
    \begin{equation}
    \label{eq:P_op_tilde}
        \bigl(\widetilde{P}^{(n)}(p)\bigr)(t) = \Bigl(\partial_t^1\,p(t)^\top,\ldots, \partial_t^n\,p(t)^\top\Bigr)^\top.
    \end{equation}
We may then deduce that any $(p, q)\in S^I_{[0,T]}$ satisfies the equation
    \begin{equation}
    \label{eq:observability_eq_tilde}
        \widetilde{F}^{(n)}\bigl(p(t), q(t) \bigr) = \bigl(\widetilde{P}^{(n)}(p)\bigr)(t),
    \end{equation}
    for all $t\in [0, T]$.
    Similarly, define operator $P^{(n)}\colon C^k([0,T];\R^{d_p})\to \bigotimes_{j=0}^{n} C^{k-j}([0,T];\R^{d_p})$ 
    by
    \begin{equation}
    \label{eq:P_op}
        \bigl(P^{(n)}(p)\bigr)(t) = \Bigl(\partial_t^0\, p(t)^\top, \partial_t^1\,p(t)^\top,\ldots, \partial_t^n\,p(t)^\top\Bigr)^\top,
    \end{equation}
    to deduce that any $(p, q)\in S^I_{[0,T]}$ satisfies the equation
    \begin{equation}
    \label{eq:observability_eq}
        F^{(n)}\bigl(p(t), q(t) \bigr) = \bigl(P^{(n)}(p)\bigr)(t),
    \end{equation}
    for all $t\in [0, T]$. 
    Notice that in the topology of \eqref{eq:norm}, $P^{(n)}$ is continuous for all $n$ since 
    differentiation $(k-j)$-times is continuous in each 
    $C^{k-j}([0, T];\R^{d_p})$ component for $0 \leq j \leq n$. 
    Similar comments apply to $\widetilde{P}^{(n)}.$

\subsection{Observability-Rank Condition}
\label{ssec:obs}





    Either of the identities \eqref{eq:observability_eq_tilde} or \eqref{eq:observability_eq}  may be
    used to define an equation for determination of $q(t)$ from knowledge of $p(\cdot)$ and its derivatives at $t.$
    We first discuss this by using $\widetilde{P}^{(n)}$, which is perhaps the most direct way to consider the
    question of recovering $q(t);$ we then modify and frame in terms of $P^{(n)}$ enabling us to connect
    with the control theory literature.
    
    Fixing some $t\in[0,T]$, and choosing 
    $\widetilde{L}:\R^{nd_p}\to \R^{d_q}$, we consider the following equation for $\fq$:
    \begin{equation}
    \label{eq:observability_eq_inv3}
    \widetilde{L}\widetilde{F}^{(n)}\bigl(p(t), \mathfrak{q} \bigr) = \widetilde{L}\bigl(\widetilde{P}^{(n)}(p)(t)\bigr).
    \end{equation}
    Given $p(\cdot)$ and its derivatives at $t$, 
    this constitutes a system of $d_q$ equations for $\fq \in \R^{d_q}$. 
     If there is a time $t$ and linear operator $\widetilde{L}$ such that the operator $\widetilde{L}\widetilde{F}^{(n)}(p(t),\,\placeholder):\R^{d_q}\to\R^{d_q}$ is invertible, then there exists a unique solution $\fq$ 
     to \eqref{eq:observability_eq_inv3} and, by \eqref{eq:observability_eq_tilde}, $\fq=q(t).$ 
     This then facilitates determination of $q(t)$ from knowledge of $p(\cdot)$ and its derivatives at $t.$
     
    % This fact provides motivation for imposing a local invertibility condition to allow existence of a map from observed $p\in C^k([0,T];\R^{d_p})$ to unobserved $q\in C^k([0,T];\R^{d_q})$ in the neighborhood of invertibility. However, in order to connect with formulations of related problems in the control theory literature, we formulate
     %a slight modification  of the foregoing.
     
     A similar approach can be undertaken using $P^{(n)}$ rather than  $\widetilde{P}^{(n)}$.
     Again fixing some $t\in[0,T]$, and now choosing 
    ${L}:\R^{(n+1)d_p}\to \R^{d_p+d_q}$, we consider the equation
    \begin{equation}
    \label{eq:observability_eq_inv2}
    {L}{F}^{(n)}\bigl(\fp, \mathfrak{q} \bigr) = {L}\bigl({P}^{(n)}(p)(t)\bigr).
    \end{equation}
    This constitutes a system of $d_p+d_q$ equations for $(\fp,\fq) \in \R^{d_p+d_q}$.
     If there exists a time $t$ and linear transformation $L$ such that the operator ${L}{F}^{(n)}(\placeholder,\,\placeholder):\R^{d_p+d_q}\to\R^{d_p+d_q}$ is invertible, then there exists a unique solution $(\fp,\fq)$, to \eqref{eq:observability_eq_inv2} and, by \eqref{eq:observability_eq}, $(\fp,\fq)=\bigl(p(t),q(t)\bigr)$ and, in particular, we recover $q(t).$ 
     This motivates the following assumption.
     
     
\begin{assumption}[Observability-Rank Condition]
\label{assump:observability}
    We say that the system \eqref{eq:dynamical_system_intro} satisfies the observability-rank condition at $(\mathfrak{p}, \mathfrak{q})\in \R^{d_p + d_q}$, if there exists $n\in\N$ and a linear operator $L\colon\R^{(n+1)d_p}\to\R^{d_p+d_q}$ such that the rank of the matrix $\sD LF^{(n)}\bigl(\fp, \fq \bigr)$ is $d_p + d_q$.
\end{assumption}


    We recall that the inverse function theorem for Euclidean spaces states that for $G\colon \R^{d} \to \R^d$ a $C^1$ map, if $\sD G$ is invertible at a point $x \in \R^d$ then there exist open neighborhoods $U, V$ of $x$ and $G(x)$ respectively, such that there exists a continuous inverse $G^{-1}\colon V \to U$. The existence of open sets $U$ and $V$ along with continuous inverse is what we refer to as local invertibility. Thus:
    
    \begin{proposition}
Assume that the  observability-rank condition holds at $(\mathfrak{p}, \mathfrak{q})\in \R^{d_p + d_q}$. Then $LF^{(n)}(\placeholder, \placeholder)$ is locally invertible at $(\fp,\fq)$.
    \end{proposition}








\begin{remark}[Observability in Control Theory]
    We note that the observability-rank condition from \Cref{assump:observability} is an instance of the observability condition from the seminal paper \cite{hermann1977nonlinear}. In particular, considering the observation function $h(\mathfrak{p},\mathfrak{q}) = \mathfrak{p}$ for any $\mathfrak{p}\in\R^{d_p}, \mathfrak{q}\in \R^{d_q}$, we may reformulate \eqref{eq:observability_eq} in terms of this observation function $h$ and hence \Cref{assump:observability} in the control theoretic setting of \cite{hermann1977nonlinear}. We elaborate on this connection in \Cref{appendix:observability_control}.
\end{remark}


\subsection{Observability of Lorenz `63 Model}
\label{ex:Lorenz}


Consider the Lorenz `63 dynamical system in the form
\begin{subequations}\label{eq:L63_dynamic}
\begin{align} 
\dot{x} &= \sigma (y - x), \\ 
\dot{y} &= x (\rho - z) - y, \\ 
\dot{z} &= xy - \beta z. 
\end{align} 
\end{subequations}
initialized at $(x_0,y_0,z_0)\in \R^3$. We consider the setting where the component $p = x$ is observed and $q = (y, z)$ is unobserved. It is readily shown that the system satisfies the observability-rank condition, given this choice
of observation. We note that by \eqref{eq:observability_eq} it holds that
\begin{align}
    &F^{(2)}(x_0, y_0, z_0)
    = \begin{pmatrix}
        x_0\\
        \sigma(y_0 - x_0)\\
        \sigma\Bigl(x_0(\rho - z_0) - y_0 - \sigma(y_0 - x_0)\Bigr)\\
    \end{pmatrix}.
\end{align}
Now, for $L$ the identity matrix, the associated differential is
\begin{align*}
    \sD\, LF^{(2)}(x_0, y_0, z_0) = \begin{pmatrix}
        1 & 0 & 0\\
        -\sigma & \sigma & 0\\
        \sigma (\rho + \sigma - z_0) & -\sigma (\sigma + 1) & -\sigma x_0
    \end{pmatrix}.
\end{align*}
From the upper-triangular structure it follows readily that this is of full rank for $x_0 \ne 0$. Hence we can conclude that the Lorenz `63 dynamical system satisfies the observability-rank condition at $t=0$, whenever $x_0\neq 0$.
In particular, the map $\{x(t)\}_{t\in [0,T]} \mapsto \{\bigl(y(t),z(t)\bigr)\}_{t\in [0,T]}$ is
well-defined whenever $x_0 \ne 0.$

To contrast with observing $\{x(t)\}_{t\in [0,T]}$, we note that the Lorenz `63 model is \textit{not} observable if the only data given is $\{z(t)\}_{t\in [0,T]}$. Specifically, if $(x,y,z)$ satisfies \eqref{eq:L63_dynamic}, then
$(-x,-y,z)$ is also a solution to \eqref{eq:L63_dynamic}. Thus, there cannot exist a well-defined mapping 
$\{z(t)\}_{t\in [0,T]} \mapsto \{\bigl(x(t),y(t)\bigr)\}_{t\in [0,T]}.$
We demonstrate this experimentally in \Cref{sec:experiments}. 


\section{Universal Approximation of Smoothing and Forecasting Maps}
\label{sec:UAforDA}


In this section we state and prove universal approximation results for neural operators approximating
the maps underlying the smoothing  \Cref{prob:smoothing} and forecasting  \Cref{prob:forecasting}.  
In \Cref{ssec:UAT} we recall a specific form of universal approximation theorem for neural operators that we employ for
both the smoothing and forecasting problems. \Cref{subsec:smoothing} is devoted to the existence and
universal approximation map $\{p(t)\}_{t\in [0,T]} \mapsto \{q(t)\}_{t\in [0,T]}$. The recovery of this map amounts to solving an instance of a smoothing \Cref{prob:smoothing}. In \Cref{subsec:forecasting} we apply the results from the subsection preceding it to deduce the existence and universal approximation of map  
$\{p(t)\}_{t\in [0,T]} \mapsto \{p(t)\}_{t\in [T,T+\tau]}$. The recovery of this map amounts to solving 
an instance of a forecasting \Cref{prob:forecasting}.

\subsection{Universal Approximation Theorem} \label{ssec:UAT}

The following proposition gives conditions under which a continuous operator may be approximated by a transformer neural operator; the result is proved in \cite[Theorem~22]{calvello2024continuum}. We use this 
specific universal approximation result in the two subsequent subsections to study the approximation of smoothing and forecasting maps. We note, however, that other neural operators could also be used \cite{kovachki2024operator}.

\begin{proposition}
\label{thm:UANeuralOperator}
Let $D \subset \R^d$ be a bounded domain with Lipschitz boundary, and fix integers $s,s'\geq 0$. If $
\Psi^\dagger\colon C^{s}(\overline{D};\,\R^{r}) \to C^{s'}(\overline{D};\,\R^{r'})$
is a continuous operator and $K \subset C^{s}(\overline{D};\,\R^{r})$ a compact set, then for any $\varepsilon>0$ there exists a transformer neural operator 
$\Psi(\placeholder; \theta)\colon K \subset C^{s}(\overline{D};\,\R^{r}) \to C^{s'}(\overline{D};\,\R^{r'})$
so that
\begin{equation}
\sup_{u\in K} \left\| \Psi^\dagger(u) - \Psi(u;\theta) \right\|_{C^{s'}} \le \varepsilon .
\end{equation}
\end{proposition}



\subsection{Smoothing}
\label{subsec:smoothing}
In this subsection we show that if a dynamical system of the form \eqref{eq:dynamical_system_intro} satisfies the regularity conditions in \Cref{assump:regularity}, together with the observability-rank condition \Cref{assump:observability}, then it is possible to construct a continuous operator mapping the observed trajectory to the initial condition of the full trajectory; we then show how this leads to existence of a continuous operator mapping the observed trajectory $p$ over a time interval to the unobserved trajectory $q$ over the same interval. We note that the continuity of these operators is necessary to apply existing universal approximation results for neural operators.

\begin{proposition}[Existence of map $W:p\mapsto (p(0), q(0))$]
\label{prop:W}
Consider the dynamical system \eqref{eq:dynamical_system_intro} satisfying the regularity \Cref{assump:regularity}, for integer $k$, and the observability \Cref{assump:observability} at some $(\mathfrak{p},\mathfrak{q})\in\R^{d_p + d_q}$, for integer $n \le k.$ Let $U \ni (\fp, \fq)$ and $V \ni LF^{(n)}(\fp, \fq)$ be the open sets given by the inverse function theorem. Then, for every compact $I \subset U$, there exists a continuous map $W\colon S^I_{[0,T],p} \subset C^{k}([0,T];\R^{d_p}) \to \R^{d_p + d_q}$ such that, for every $(p, q) \in S^I_{[0,T]}$,
\begin{equation}
\label{eq:W}
W(p) = (p(0), q(0)).
\end{equation}
\end{proposition}
\begin{proof}
Observe that from \eqref{eq:observability_eq_inv2} for every $(p, q) \in S^I_{[0,T]}$ we have that 
\begin{equation}
\label{eq:W-capability}
    LF^{(n)}(p(0), q(0)) = L(P^{(n)}(p)(0))
\end{equation}
and $LF^{(n)}(\placeholder, \placeholder)$ is locally invertible by the inverse function theorem from \Cref{assump:observability}. Furthermore, the evaluation functional at time $t=0$, namely $\delta_0 \colon \bigotimes_{j=0}^n C^{k-j}([0,T];\R^{d_p}) \to \R^{(n+1)d_p}$, is linear and bounded, hence continuous with respect to the topology of \eqref{eq:norm}. Using these operators, we can construct $W\colon S^I_{[0,T],p} \subset C^{k}([0,T];\R^{d_p}) \to \R^{d_p + d_q}$ so that
\begin{equation}
      W(p) = \bigl(LF^{(n)}\bigr)^{-1}
      \circ L \circ \delta_0 \circ P^{(n)}(p).
\end{equation}
From \eqref{eq:W-capability} we have that $ L \circ \delta_0 \circ P^{(n)}(p) \in V$ for all $p$, so composition with $(LF^{(n)})^{-1}$ is well-defined. By construction $(LF^{(n)})^{-1}$ is continuous and $W$ is continuous since it is a composition of continuous operators. For all $(p, q) \in S^I_{[0,T]}$ we have $W(p) = (p(0), q(0))$ since $(p(0), q(0))$ is a solution to \eqref{eq:W-capability} and $LF^{(n)}$ is a bijection from $U$ to $V$. 
\end{proof}

Following \Cref{prop:W} we may now establish the existence of an operator mapping observed $p$ trajectories to unobserved $q$ trajectories.
\begin{proposition}[Existence of map $\Psi_S^\dagger:p\mapsto q$]
\label{prop:psiS}
Consider the dynamical system \eqref{eq:dynamical_system_intro} satisfying
regularity \Cref{assump:regularity}, for integer $k$, and the observability \Cref{assump:observability} at some $(\mathfrak{p},\mathfrak{q})\in\R^{d_p + d_q}$, for integer $n \le k.$
Let $U \ni (\fp, \fq)$ and $V \ni LF^{(n)}(\fp, \fq)$ be the open sets given by the inverse function theorem. Then, for every compact $I \subset U$, there exists continuous operator $\Psi_S^\dagger \colon S^I_{[0,T],p} \subset C^k([0,T];\R^{d_p}) \to S^I_{[0,T],q} \subset C^k([0, T]; \R^{d_q})$ such that, for every $(p, q) \in S^I_{[0,T]}$
and $t \in [0,T]$,
\begin{equation}
    \label{eq:p_to_q_map1}
    \Psi_S^\dagger(p)(t) = q(t).
\end{equation}
\end{proposition}
\begin{proof}
We construct such a $\Psi^\dagger$ using $W$ from \Cref{prop:W} and composition with the solution operator $\Phi$. Define $\Psi_S^\dagger : S^I_{[0,T],p} \subset C^k([0,T];\R^{d_p}) \to S^I_{[0,T],q} \subset C([0, T]; \R^{d_q})$ by
\begin{equation}
    \label{eq:p_to_q_map2}
    \Psi_S^\dagger(p)(t) = \pi_q \circ \Phi(t, W(p)), \quad 0\leq t\leq T.
\end{equation}
Recall that $\pi_q$ is continuous. By construction $W(p) = (p(0), q(0))$, so $\pi_q \circ \Phi(t, W(p)) = q(t)$. Observe that this is a continuous operator since $\Phi$ is continuous with respect to initial conditions, $W$ is continuous with respect to $p$ and composition of continuous operators is continuous.
\end{proof}
Leveraging the approximation properties of neural operators, we may now establish a universal approximation theorem for the solution of \Cref{prob:smoothing}, an instance of a smoothing problem. In the following, set $I$ is defined through the two preceding propositions.

\begin{theorem}[Universal Approximation for Smoothing]
\label{thm:UAsmoothing}
Consider the dynamical system \eqref{eq:dynamical_system_intro} satisfying 
regularity \Cref{assump:regularity}, for integer $k$, and the observability \Cref{assump:observability} at some $(\mathfrak{p},\mathfrak{q})\in\R^{d_p + d_q}$, for integer $n \le k.$
For any $\epsilon > 0$ there exists a transformer neural operator $\Psi(\placeholder;\theta): S^I_{[0,T],p} \subset C^k ([0,T];\R^{d_p}) \to  S^I_{[0,T],q} \subset C^k ([0,T];\R^{d_q})$ satisfying
\[\sup_{p \in S^I_{[0,T],p}} \|\Psi_S^\dagger(p) - \Psi(p;\theta)\|_{C^k} \leq \epsilon. \]
\end{theorem}

\begin{proof}
We know that $\Psi_S^{\dagger}$ exists and continuous since the assumptions of \Cref{prop:psiS} are satisfied. We recall that $S^I_{[0,T],p}$ is compact. The conclusion follows from a direct application of \Cref{thm:UANeuralOperator} with $D = [0, T]$, $r = d_p$, $r' = d_q$, $\Psi^{\dagger} = \Psi_S^{\dagger}$, and $K = S_{[0,T],p}$.
\end{proof}

\begin{remark}
\label{rem:global}
    Since our observability-rank condition \Cref{assump:observability} is local, all operators constructed in the previous theorems have a domain of input trajectories which are near the point of invertibility $(\fp,\fq)$. The set of these trajectories is determined, non-constructively, by the inverse function theorem. It is of general interest to determine the largest possible domain of input trajectories for which the smoothing map $\Psi^\dagger_S$ (or the subsequently defined forecasting map $\Psi^\dagger_F$) exists and is continuous. A potential approach to this problem is to determine all points of invertibility and stitch together a global inverse from the open sets given by the inverse function theorem. Such analysis, however, will likely need to be carried out on a case-by-case basis and is not the focus of the current work.

    
    In the setting of the Lorenz `63 system studied in \Cref{ex:Lorenz}, we show that $U = \mathbb{R}\setminus \{0\} \cup \mathbb{R}^2$. Therefore the domain $S^I_{[0,T],p}$ of $\Psi^\dagger_S$ contains all trajectories with initial condition $x_0$ uniformly bounded away from $0$.
\end{remark}

\subsection{Forecasting}
\label{subsec:forecasting}
Using \Cref{prop:W} we may also establish the existence of an operator mapping the observed portion of $p$ trajectories, $p \!\restriction_{[0,T]}$, to the future portion $p \!\restriction_{[T,T+\tau]}$.
\begin{proposition}[Existence of map $\Psi_F^\dagger:p \!\restriction_{[0,T]}\mapsto p \!\restriction_{[T,T+\tau]}$]
\label{prop:psiF}
Consider the dynamical system \eqref{eq:dynamical_system_intro} satisfying 
regularity \Cref{assump:regularity}, for integer $k$, and the observability \Cref{assump:observability} at some $(\mathfrak{p},\mathfrak{q})\in\R^{d_p + d_q}$, for integer $n \le k.$
Let $ U \ni (\fp, \fq)$ and $V \ni LF^{(n)}(\fp, \fq)$ be the open sets given by the inverse function theorem. 
Then, for every compact $I \subset U$, there exists a continuous operator $\Psi_F^\dagger \colon S^I_{[0,T],p} \subset C^k([0,T];\R^{d_p}) \to S^I_{[T,T+\tau],p} \subset C^k([T, T+\tau]; \R^{d_p})$ such that, for every $(p, q) \in S^I_{[0,T]}$,
\begin{equation}
    \label{eq:p_to_q_map3}
    \Psi_F^\dagger(p)(t) = p(t), \qquad T\leq t\leq T+\tau.
\end{equation}
\end{proposition}
\begin{proof}
We construct such a $\Psi_F^\dagger$ using $W$ from \Cref{prop:W} and composition with the solution operator $\Phi$. Define $\Psi_F^\dagger \colon S^I_{[0,T],p} \subset C^k([0,T];\R^{d_p}) \to S^I_{[T,T+\tau],p} \subset C^k([T, T+\tau]; \R^{d_p})$ by 
\begin{equation}
    \label{eq:p_to_q_map4}
    \Psi_F^\dagger(p)(t) = \pi_p \circ \Phi(t, W(p)), \quad T\leq t\leq T+\tau.
\end{equation}
We first recall that $\pi_p$ is continuous. By construction it holds that $W(p) = (p(0), q(0))$, hence $\pi_p \circ \Phi(t, W(p)) = {p}(t)$ for $T\leq t\leq T+\tau$. We observe that this is a continuous operator since $\Phi$ is continuous with respect to initial conditions, $W$ is continuous with respect to $p$ and composition of continuous operators is continuous.
\end{proof}
Leveraging the approximation properties of neural operators, we may now establish a universal approximation theorem for the solution of \Cref{prob:forecasting}, an instance of a forecasting problem.

\begin{theorem}[Universal Approximation for Forecasting]
\label{thm:UAforecasting}
Consider the dynamical system \eqref{eq:dynamical_system_intro} satisfying
regularity \Cref{assump:regularity}, for integer $k$, and the observability \Cref{assump:observability} at some $(\mathfrak{p},\mathfrak{q})\in\R^{d_p + d_q}$, for integer $n \le k.$ For any $\epsilon > 0$ there exists a transformer neural operator $\Psi(\placeholder;\theta)\colon S^I_{[0,T],p} \subset C^k ([0,T];\R^{d_p}) \to S^I_{[T,T+\tau],p} \subset C^k ([T,T+\tau];\R^{d_p})$ satisfying
\[\sup_{p \in S^I_{[0,T],p}} \|\Psi_F^\dagger(p) - \Psi(p;\theta)\|_{C^k([T,T+\tau];\R^{d_p})} \leq \epsilon. \]
\end{theorem}
\begin{proof}
    We know that $\Psi_F^{\dagger}$ exists and is continuous since the assumptions of \Cref{prop:psiF} are satisfied. Recall that $S^I_{[0,T],p}$ is compact. We apply the result of \Cref{thm:UANeuralOperator} to an operator defined via linear homeomorphisms of $\Psi_F$, which we will then use to establish the approximation result in its final form. Indeed, we define the translation map $\Lambda_T:C^k([T,T+\tau];\R^{d_p})\to C^k([0,\tau];\R^{d_p})$ so that
    \[(\Lambda_T \varphi)(s) = \varphi(T+s), \qquad s\in[0,\tau], \]
    for any $\varphi\in C^k([T,T+\tau];\R^{d_p})$. We note that 
    \[\|\varphi\|_{C^k([T,T+\tau];\R^{d_p})} = \|\Lambda_T\varphi\|_{C^k([0,\tau];\R^{d_p})}.\]
    It is also possible to define the rescaling operator $R_\tau:C^k([0,\tau];\R^{d_p})\to C^k([0,1];\R^{d_p})$ so that
    \[
    (R_\tau\varphi)(s) = \varphi(\tau s), \qquad s\in[0,1],
    \]
    for any $\varphi\in C^k([0,\tau];\R^{d_p})$. We similarly define the rescaling operator $R_T$. We note that these rescaling operators are linear homeomorphisms, preserving equivalence of norms up to multiplicative constants. We therefore define the translated and rescaled forecasting operator $\widetilde{\Psi}^\dagger_F:R_T(S^I_{[0,T],p})\subset C^k([0,1];\R^{d_p})\to C^k([0,1];\R^{d_p})$ defined via the composition
    \[
    \widetilde{\Psi}^\dagger_F = R_\tau \circ \Lambda_T \circ {\Psi}^\dagger_F \circ R_T^{-1},
    \]
    which is continuous by continuity of all the operators in the composition. We note that $R_T(S^I_{[0,T],p})$ is compact since $R_T$ is a homeomorphism and therefore preserves compactness. Hence, by \Cref{thm:UANeuralOperator}, it is possible to deduce that for any $\delta>0$ there exists a neural operator $\widetilde{\Psi}(\placeholder,\theta): R_T(S^I_{[0,T],p})\subset C^k([0,1];\R^{d_p})\to C^k([0,1];\R^{d_p})$ satisfying
    \begin{equation}
    \label{eq:approximant}
    \sup_{p \in R_T(S^I_{[0,T],p})} \|\widetilde{\Psi}_F^\dagger(p) - \widetilde{\Psi}(p;\theta)\|_{C^k([0,1];\R^{d_p})} \leq \delta. 
    \end{equation}
    This indeed follows from a direct application of \Cref{thm:UANeuralOperator} with $D = [0, 1]$, $r = d_p$, $r' = d_p$, $\Psi^{\dagger} = \widetilde{\Psi}_F^{\dagger}$, and $K = R_T(S^I_{[0,T],p})$. Defining $\Psi(\placeholder,\theta):C^k ([0,T];\R^{d_p}) \to C^k ([T,T+\tau];\R^{d_p})$ via the composition
    \[
    \Psi(\placeholder,\theta) = \Lambda_T^{-1}\circ R_\tau^{-1} \circ \widetilde{\Psi}(\placeholder;\theta)\circ R_T,
    \]
    it is readily deduced that for any $p \in S^I_{[0,T],p}$ it holds that
    \[
    \|\Psi_F^\dagger(p) - \Psi(p;\theta)\|_{C^k([T,T+\tau];\R^{d_p})} = C\|\widetilde{\Psi}_F^\dagger(R_Tp) - \widetilde{\Psi}(R_Tp;\theta)\|_{C^k([0,1];\R^{d_p})},
    \]
    for a constant $C$ depending on $T$ and $\tau$. Therefore, choosing $\epsilon=\delta/C$ and applying the result established in \eqref{eq:approximant} yields the conclusion.
    \end{proof}

\section{Data-Driven Approximation of Smoothing and Forecasting Maps}
\label{sec:experiments}

The theoretical developments of \Cref{sec:UAforDA}  prove the existence of, and universal approximation theorems for,
maps that underpin data-driven smoothing and forecasting. In this section we demonstrate that
it is possible to learn these maps, purely from data, in practice. In \Cref{ssec:esu} we describe the experimental
set-up in general, followed in \Cref{ssec:summ} by a summary of the numerical results.
We describe details of the experiment with the Lorenz `63 \cite{lorenz63} and Lorenz `96 \cite{lorenz96} dynamical systems in \Cref{subsec:lorenz63} and \Cref{subsec:lorenz96}, respectively; experiments with the Kuramoto-Sivashinsky equation \cite{kuramoto1978diffusion, michelson1977nonlinear} are contained in \Cref{subsec:lorenz96}.


\subsection{Experimental Set-Up} \label{ssec:esu} To train the neural operators solving the data assimilation problems we consider the following data scenario. For each dynamical system, the set $\{p^{(j)}, q^{(j)} \}_{j=1}^J$ of input-output pairs are generated using the dynamics in \eqref{eq:dynamical_system_intro}, given i.i.d. initial conditions $\bigl(p^{(j)}(0),q^{(j)}(0)\bigr)\sim\nu$. Measure $\nu$ is computed as the pushforward, over some burn-in time, of the simpler measure $\nu_0.$ Burn-in is used to ensure that $\nu$ is close to the invariant measure
supported on the global attractor of the dynamical system. For each dynamical system we consider the existence of operators $\Psi^\dagger_{S}:C([0,T])\to C([0,T])$ and $\Psi^\dagger_{F}:C([0,T])\to C([T,T+\tau])$ defined as the mappings\footnote{We note that the existence of the operators in \eqref{eq:ops_numerics} is not directly guaranteed by \Cref{prop:psiS} and \Cref{prop:psiF}, as the former are defined on the whole of $C([0,T])$, rather than a local domain of existence. We note that establishing the existence of the operators in \eqref{eq:ops_numerics} constitutes an important avenue for future work, which we highlight in \Cref{sec:conclusions}. We assume their existence for our numerical investigations.}
\begin{equation}
\label{eq:ops_numerics}
    \Psi^\dagger_{S}:\{p(t)\}_{t\in [0,T]}\mapsto \{q(t)\}_{t\in [0,T]}\quad \Psi^\dagger_{F}:\{p(t)\}_{t\in [0,T]}\mapsto \{p(t)\}_{t\in [T,T+\tau]}.
\end{equation}
We train transformer neural operators \cite{calvello2024continuum} to construct approximations of $\Psi^\dagger_{S}$ and $\Psi^\dagger_{F}$ picked from the parametric class
$$\Psi_{S}:C([0,T])\times \Theta \to C([0,T]), \quad \Psi_{F}:C([0,T])\times \Theta \to C([T,T+\tau])$$
respectively. In all cases we use the transformer neural operator based on self-attention from \cite{calvello2024continuum} to approximate the $\Psi^\dagger_{S}$. To approximate $\Psi^\dagger_{F}$ in the context of Lorenz `63 dynamics we also use the same self-attention based transformer neural operator. However, to approximate 
$\Psi^\dagger_{F}$ for the Lorenz `96 and Kuramoto-Sivashinsky models we employ a variant of the transformer neural operator architecture which uses cross-attention. Architectural details as well as the universal approximation properties for these neural operators are discussed in more depth in \Cref{appendix:architectures}. 

The set $\Theta$ is assumed to be a finite dimensional parameter space from which a $\theta^*\in\Theta$ is selected so that $\Psi_{\placeholderB}(\placeholder,\theta^*) \approx \Psi_{\placeholderB}^\dagger$. (Here $\placeholderB$ is either $S$ or $F$.) To define
$\theta^*$ we employ a cost functional $c$ and solve the following optimization problem:
\begin{equation}
\label{eq:data_assimilation_model}
\theta^*={\arg \min}_{\theta\in\Theta}\mathbb{E}_{\placeholder}\Bigl[ c\bigl(\Psi_{\placeholderB}(\placeholder,\theta), \Psi_{\placeholderB}^\dagger(\placeholder) \bigr)\Bigr].
\end{equation}
In practice, we will only have access to the values of each input and output trajectory at a set of discretization points of the domain, which we denote as $\{t_i\}_{i=1}^N$. The self-attention based transformer neural operator that we deploy for smoothing uses the same number of grid-points for the input and output; employing a cross-attention based transformer neural operator is used to bypass this constraint on matching input and output grids for forecasting in some settings, where the input is on the time-interval $[0,T]$ and the output on $[T,T+\tau]$. We note that due to its discretization invariant properties, this architecture can be used to obtain a prediction for the output at any arbitrary discretization of $[T,T+\tau]$. For the smoothing experiments we use $6$ self-attention transformer neural operator layers with $128$ latent channels, for each self-attention operator, we use $8$ attention heads. We also use this set-up for forecasting in the Lorenz `63 context. For the other forecasting experiments we use an architecture composed of an encoder and a decoder: the encoder is comprised of $6$ self-attention transformer neural operator layers with $1024$ latent channels; the decoder on the other hand is comprised of a single layer consisting of a self-attention operator and a cross-attention operator both with $1024$ latent channels. 

\subsection{Summary of Results from Numerical Experiments} \label{ssec:summ}

We provide here a summary of the experimental results. Details of the exact experimental set-up, for each problem, are
considered in the subsections that follow. We start, in \Cref{tab:smoothing_results}, where
we collect the results of the smoothing experiments for the various systems, reporting relative $L^2$ errors between the true unobserved trajectory and the prediction generated by the neural operator. All errors are reported on test data not used for training. 
 
\begin{table}[h]
\footnotesize
\centering
\caption{Smoothing performance on Lorenz `63 (L63), Lorenz `96 (L96), and the Kuramoto--Sivashinsky equation (KSE).}
\label{tab:smoothing_results}
\begin{tabular}{lccccccc}
\toprule
\textbf{System} 
& \textbf{Avg. Rel. $L_2$} 
& \textbf{Med. Rel. $L_2$} 
& \textbf{Std. Rel. $L_2$} 
& \textbf{Min. Rel. $L_2$} 
& \textbf{Max. Rel. $L_2$} 
%& \textbf{Mean Err} 
%& \textbf{Var Err} 
\\
\midrule
L63 
& 0.012426 
& 0.012402 
& 0.001612 
& 0.008761 
& 0.017275 
%& 0.003013 
%& 0.007778 
\\
L96 
& 0.047156 
& 0.044635 
& 0.014569 
& 0.025633 
& 0.258743 
%& 0.011351 
%& 0.019663 
\\
KSE 
& 0.009433 
& 0.009322 
& 0.000755 
& 0.007678 
& 0.013157 
%& 0.001073 
%& 0.004319 
\\
\bottomrule
\end{tabular}
\end{table}

In \Cref{tab:forecasting_results} we report the results of the forecasting experiments. For each system, we report the mean, median and standard deviation of the relative $L^2$ errors obtained on the sample trajectories in the test set. We also report the relative improvement achieved by the neural operator prediction when compared to the constant prediction taken using the final value of the input trajectory. We note that the chaotic nature of the problems considered means that long-term trajectory forecasting is difficult. However we can compose forecasts and look
at the resulting statistics of the time-series -- in particular the invariant measure of the resulting stationary process; these are well-approximated. Thus, for each of the systems, we report in the respective subsections the statistics of the long-time forecasts obtained via composition of the trained neural operators. The results in \Cref{tab:forecasting_results} along with the accurate prediction of trajectory statistics demonstrate the success of using a neural operator forecasting approach.

\begin{table}[htbp]
\footnotesize
\centering
\caption{Forecasting performance on Lorenz `63 (L63), Lorenz `96 (L96), and the Kuramoto--Sivashinsky equation (KSE). In addition to the relative $L^2$ errors achieved on the test set, we also report the relative improvement of $L^2$ errors achieved with the transformer neural operator compared to a constant forecast equal to the last input state.}
\label{tab:forecasting_results}
\begin{tabular}{lcccc}
\toprule
\textbf{System} 
& \textbf{Avg. Rel. $L_2$} 
& \textbf{Med. Rel. $L_2$} 
& \textbf{Std. Rel. $L_2$} 
& \textbf{Rel. Improvement (\%)} \\
\midrule
L63  & 0.050 & 0.033 & 0.100 & 95.53 \\
L96  & 0.033 & 0.032 & 0.004 & 94.38 \\
KSE  & 0.046 & 0.045 & 0.007 & 82.68 \\
\bottomrule
\end{tabular}
\end{table}

% \begin{figure}[t]
%   \centering

%   % --- top row ---
%   \begin{subfigure}[t]{0.48\textwidth}
%     \centering
%     \includegraphics[width=\linewidth]{figures/L63_histogram_comparison.png}
%     \caption{}
%     \label{fig:tri-a}
%   \end{subfigure}\hfill
%   \begin{subfigure}[t]{0.48\textwidth}
%     \centering
%     \includegraphics[width=\linewidth]{figures/L96_histogram_comparison.png}
%     \caption{}
%     \label{fig:tri-b}
%   \end{subfigure}
%   % --- bottom centered ---
%   \begin{subfigure}[t]{0.48\textwidth}
%     \centering
%     \includegraphics[width=\linewidth]{figures/L96_histogram_comparison.png}
%     \caption{}
%     \label{fig:tri-c}
%   \end{subfigure}
%   \caption{In this panel we display the distributions of points in the predicted and ground truth trajectories. In (a) we display the result
%   \label{fig:tri}
% \end{figure}

% \begin{table}[htbp]
% \footnotesize
% \centering
% \caption{Comparison of trajectory-level error statistics between learned forecasting models and constant baselines across L63, L96, and KSE systems.}
% \label{tab:trajectory_stats}
% \begin{tabular}{lcccc}
% \toprule
% \textbf{Errors} 
% & \textbf{Mean (Learned)} 
% & \textbf{Mean (Baseline)} 
% & \textbf{Var. (Learned)} 
% & \textbf{Var. (Baseline)} \\
% \midrule
% L63 & 0.034 & 0.215 & 0.0023 & 0.0238 \\
% L96 & 0.048 & 0.187 & 0.0079 & 0.0217 \\
% KSE & 0.032 & 0.093 & 0.0026 & 0.0055 \\
% \bottomrule
% \end{tabular}
% \end{table}

In the following subsections, we describe the experimental set-up for the various systems and provide analysis of the numerical results.

\subsection{Lorenz `63}
\label{subsec:lorenz63}

We study the Lorenz `63 dynamical system \cite{lorenz63} in the form given in \eqref{eq:L63_dynamic}.
For the purposes of our experiments, the parameters are set to the classical values
$\sigma = 10,\ \rho = 28,\ \beta = \tfrac{8}{3}$ at which the system has provable stable chaotic and statistical properties \cite{tucker1999lorenz,holland2007central}. The initial conditions at time $t=-20$, for both training and test sets, are sampled i.i.d from the uniform distributions $U([-15,15])$, $U([-15,15])$ and $U([0,40])$, respectively; this defines $\nu_0.$ We use a burn-in time of $20$, ensuring that the trajectories are close to the global attractor at time $t=0$, sampled from measure $\nu$, approximating the invariant measure supported on the global attractor. For the \Cref{prob:smoothing} smoothing experiment we use the $x$ trajectory on a time interval of $[0,5]$ as input data and predict the coupled $y,z$ trajectories over the same time interval. For practical implementation, we use points on the trajectories sampled at uniform intervals $\Delta t= 0.02$. For the \Cref{prob:forecasting} forecasting experiment we use the $x$ trajectory on a time interval of $[0,2]$ as input data and predict the coupled $x$ trajectory over the future time $[2,4]$. For practical implementation, we use points on the input and output trajectories sampled at uniform intervals $\Delta t= 0.01$. We train the neural operator solving the smoothing problem using $10000$ training trajectories and the one solving the forecasting problem using $200000$ trajectories. For both experiments we use a test set of $2000$ trajectories. We note the higher data complexity required to solve the forecasting problem. 



We summarize the parameter and experimental settings for the Lorenz `63 dynamical system in \Cref{tab:lorenz63_settings}.
\begin{table}[h]
\footnotesize
\centering
\caption{Lorenz `63 system settings}
\label{tab:lorenz63_settings}
\begin{tabular}{lcc}
    \toprule
    \textbf{Category} & \textbf{Smoothing} & \textbf{Forecasting} \\
    \midrule
    Parameters 
        &\multicolumn{2}{c}{ $\sigma = 10,\ \rho = 28,\ \beta = \tfrac{8}{3}$} \\
    \midrule
    Observed 
        & $x \in C\bigl([0,5]; \bbR \bigr)$ 
        & $x \in C\bigl([0,2]; \bbR \bigr)$ \\
    \midrule
    Unobserved 
        & $(y,z) \in C\bigl([0,5]; \bbR^2 \bigr)$ 
        & $(y,z) \in C\bigl([2,4]; \bbR^2 \bigr)$ \\
    \midrule
    Observation time step 
        & $\Delta t = 0.02$
        & $\Delta t = 0.01$ \\
    \bottomrule
\end{tabular}
\end{table}
In \Cref{fig:L63_smoothing} we demonstrate the qualitative results of the smoothing experiment. Indeed, we display the predicted trajectories corresponding to the median and highest relative $L^2$ errors in the test set compared to true trajectories. The panels showcase predictions that are qualitatively indistinguishable from the truth, indicating the success of the purely data-driven smoothing approach.
\begin{figure}
    \centering
    \includegraphics[width=\linewidth]{figures/L63_smoothing_median_worst_comparison.png}
    \caption{Lorenz `63. Median and worst-case relative $L^2$ error samples for smoothing experiment involving prediction of $(y,z)$ from $x$.}
    \label{fig:L63_smoothing}
\end{figure}

To highlight the importance of observability, we experiment with the smoothing setting of mapping the $z$ trajectory over the time interval $[0,5]$ to the $x,y$ trajectories over the same time interval. We use the same model and training setup as before. \Cref{fig:L63_smoothing_xy} displays the predicted trajectories corresponding to the median and highest relative $L^2$ errors in the test set compared to true trajectories. The panels showcase the failure in the predictions, which achieve a median relative $L^2$ error of $1$; it is also interesting to note that the model seems to predict the $0$ trajectory, possibly due to the reflection argument discussed in \Cref{ex:Lorenz}.
\begin{figure}
    \centering
    \includegraphics[width=\linewidth]{figures/L63_median_worst_comparison_xy_pred.png}
    \caption{Lorenz `63. Median and worst-case relative $L^2$ error samples for smoothing experiment involving prediction of $(x,y)$ from $z$.}
    \label{fig:L63_smoothing_xy}
\end{figure}

In \Cref{fig:L63_forecasting_all} we demonstrate the qualitative results of the forecasting experiment, based on data from the $x$ coordinate. Indeed, in \Cref{fig:L63_forecasting} we display the predictions for the trajectories corresponding to the median and highest relative $L^2$ errors in the test set compared to ground truths. We note that the prediction for the test sample presenting the largest error still resembles a possible trajectory of the Lorenz `63 system. Indeed, the prediction diverges to another lobe of the attractor; this behavior is to be expected in prediction problems in the context of chaotic systems. With the goal of obtaining forecasts on longer time horizons, we experiment with composition of the learned forecasting map $\Psi_F:C([0,2])\to C([2,4])$ to obtain $\Psi_F^n:C([0,2])\to C([2,2+2n])$. In \Cref{fig:L63_forecasting_hist}, for $n=500$ we compare the distribution (histogram) of points in the predicted future trajectories $\widehat{p}$ to points in the ground truth trajectories $p$. The matching of distributions, and hence statistics, indicates the success in predicting valid trajectories from the attractor of the Lorenz system, despite the chaotic nature of the equations and hence per-trajectory divergence. As further validation of the quantitative success of the purely data-driven forecasting approach, we recall the numerical comparison with the prediction given by the constant value of $x(T)$ for $T=2$, that is taking the predicted trajectory to be $\widehat{p}(t)=x(2)$ for $t\in[2,4]$. \Cref{tab:forecasting_results} shows the improvement in accuracy yielded by the cross-attention based transformer neural operator, in comparison to this naive approach.
\begin{figure}[htbp]
    \centering
    \begin{subfigure}{0.48\linewidth}
        \centering
        \includegraphics[width=\linewidth]{figures/L63_median_worst_comparison_2.png}
        \caption{}
        \label{fig:L63_forecasting}
    \end{subfigure}\hfill
    \begin{subfigure}{0.48\linewidth}
        \centering
        \includegraphics[width=\linewidth]{figures/L63_histogram_comparison.png}
        \caption{}
        \label{fig:L63_forecasting_hist}
    \end{subfigure}
    \caption{Lorenz `63. Left: median and worst-case relative $L^2$ error samples in forecasting.
    Right: distribution of trajectory predictions under composition of the learned forecasting map; the dark shade of red represents overlap between the blue ground truth distribution and the light red predicted distribution. The prevalence of the dark shade of red thus indicates matching of distributions.}
    \label{fig:L63_forecasting_all}
\end{figure}

\subsection{Lorenz `96}
\label{subsec:lorenz96}

The Lorenz `96 model \cite{lorenz96} consists of a linearly damped and externally forced system equipped with an energy-conserving quadratic nonlinearity, which generates cyclic interactions among the state variables. These properties make the model a standard testbed for investigating atmospheric predictability. The system dynamics are governed by the following equations:
\begin{equation}\label{eq:L96_equations}
    \frac{du_i}{dt} = \bigr(u_{i+1} - u_{i-2}\bigl)u_{i-1} - u_i + F,\quad i=1,2,\ldots, d.
\end{equation}
Periodicity is imposed to define variables $u_j$ with indices outside the set $j=1,2,\ldots, d.$
For the purposes of our experiments we set $d=40$ and $F=8$. For all components $u_i$ the initial conditions for training and test sets are found by first selecting initial conditions drawn from $F+U([-1,1])$, 
i.i.d. in both index $i$ and random choice of initialization at time $t=-200$, and  then simulating the system using a burn-in time of $200$ to ensure the initializations of the trajectories at $t=0$ are sampled i.i.d. from a measure approximating the invariant measure supported on the global attractor at time $t=0$. For the smoothing experiment we take trajectories over a time interval of $[0,5]$ as input data and predict the unobserved trajectories over the same time interval; we implement using points on the trajectories sampled at uniform intervals $\Delta t= 0.02$. For the forecasting experiment we also use trajectories on a time interval of $[0,5]$ as input data and predict observed trajectories over the future time $[5,5.2]$. We implement using points on the input and output trajectories sampled at uniform intervals $\Delta t= 0.02$. We train the neural operator solving the smoothing problem using $10000$ training trajectories and the one solving the forecasting problem using $160000$ trajectories. For both experiments we use a test set of $2000$ trajectories. We again note the higher data complexity required to solve the forecasting problem. 

We summarize the parameter and experimental settings for the Lorenz `96 dynamical system in \Cref{tab:lorenz96_settings}.
\begin{table}[h]
\footnotesize
\centering
\caption{Lorenz `96 system settings}
\label{tab:lorenz96_settings}
\begin{tabular}{lcc}
    \toprule
    \textbf{Category} & \textbf{Smoothing} & \textbf{Forecasting} \\
    \midrule
    Parameters
       &\multicolumn{2}{c}{$F=8$} \\
    \midrule
    Observed
        & $(u_1,\ldots,u_{21},u_{23},u_{25},\ldots,u_{39}) \in C([0,5];\mathbb{R}^{30})$ & $(u_1,u_{2},u_{3},u_{5},\ldots,u_{39}) \in C([0,5];\mathbb{R}^{30})$
          \\
    \midrule
    Unobserved 
        & $(u_{22},u_{24},\ldots,u_{40}) \in C([0,5];\mathbb{R}^{10})$
        & $(u_1,u_{2},u_{3},u_{5},\ldots,u_{39}) \in C([5,5.2];\mathbb{R}^{30})$ \\
    \midrule
    Obs. time
        &\multicolumn{2}{c}{ $\Delta t = 0.02$ }
        \\
    \bottomrule
\end{tabular}
\end{table}
In \Cref{fig:L96_smoothing} we demonstrate the qualitative results of the smoothing experiment. Indeed, we display the predicted trajectories corresponding to the median and highest relative $L^2$ errors in the test set compared to true trajectories. The median error panel demonstrates near-perfect overlap between prediction and truth, while the highest error panel exhibits overlap only after $t=1$, indicating a possible effect of synchronization \cite{pecora1990synchronization}. Together, these results demonstrate the qualitative success of the purely data-driven smoothing approach. 
\begin{figure}[htbp]
    \centering
    \includegraphics[width=0.5\linewidth]{figures/L96_smoothing_median_worst_comparison.png}
    \caption{Lorenz `96. Median and worst-case performance on the Lorenz `96 system for smoothing.}
    \label{fig:L96_smoothing}
\end{figure}
In \Cref{fig:L96_forecast_spatiotemporal} we display results for the forecasting experiment. Here we show the pointwise errors relative to the root mean square of the ground truth (computed in space) for the sample in the test set achieving the median relative $L^2$ error. With the goal of obtaining forecasts on longer time horizons, we experiment with composition of the learned forecasting map $\Psi_F:C([0,5])\to C([5,5.2])$ to obtain $\Psi_F^n:C([0,5])\to C([5,5+0.2n])$. In \Cref{fig:L96_forecasting}, we show an example of predicted trajectories obtained by composing the map with $n=50$; the plot shows trajectory divergence after a handful of compositions, due to chaos and error compounding. However, in \Cref{fig:L96_forecasting_hist}, for $n=500$ we compare the distribution of points in the predicted future trajectories $\widehat{p}$ to the distribution of
points in the ground truth trajectories $p$. The matching of distributions, and hence statistics, indicates the success in predicting the invariant measure supported on the attractor of the Lorenz `96 system, despite the chaotic nature of the equations that results in trajectory divergence. As further validation of the quantitative success of the purely data-driven forecasting approach, we recall the numerical comparison with the prediction given by the constant value of $x(T)$ for $T=5$, that is taking the predicted trajectory to be $\widehat{p}(t)=p(5)$ for $t\in[5,5.2]$. \Cref{tab:forecasting_results} shows the improvement in accuracy yielded by the cross-attention based transformer neural operator, in comparison to this naive approach.
\begin{figure}[htbp]
    \centering
    \includegraphics[width=0.6\linewidth]{figures/L96_spatiotemporal_error.png}
    \caption{Lorenz `96. Pointwise errors relative to the root mean square of the ground truth (computed in space) for the sample in the test set achieving the median relative $L^2$ error.}
    \label{fig:L96_forecast_spatiotemporal}
\end{figure}
\begin{figure}[htbp]
    \centering
    \begin{subfigure}{0.48\linewidth}
        \centering
        \includegraphics[width=\linewidth]{figures/L96_forecasting_median_worst_comparison.png}
        \caption{}
        \label{fig:L96_forecasting}
    \end{subfigure}\hfill
    \begin{subfigure}{0.48\linewidth}
        \centering
        \includegraphics[width=\linewidth]{figures/L96_histogram_comparison.png}
        \caption{}
        \label{fig:L96_forecasting_hist}
    \end{subfigure}
    \caption{Lorenz `96. Left: forecast obtained by composing the learned map on one-step median and worst-case relative $L^2$ error samples. Right: distribution of trajectory predictions under composition of the learned forecasting map; the dark shade of red represents overlap between the blue ground truth distribution and the light red predicted distribution. The prevalence of the dark shade of red thus indicates matching of distributions.}
    \label{fig:L96_comparison}
\end{figure}

\subsection{Kuramoto-Sivashinsky}
\label{subsec:ks}

The Kuramoto–Sivashinsky (KS) equation \cite{kuramoto1978diffusion} is a well-studied nonlinear partial differential equation that models the development of instabilities in spatially extended systems. It exhibits rich spatiotemporal chaotic behavior and arises in a variety of physical contexts, including flame front propagation, thin film flows, and reaction–diffusion phenomena. With the periodicity in space imposed to identify the solution at $x=L$ and $x=0$, the one-dimensional KS equation for $u: [0,L] \times \mathbb{R}^+ \to \mathbb{R}$ takes the form
\begin{subequations}\label{eq:KS_equation}
\begin{align}
    \frac{\partial u}{\partial t} + \frac{\partial^4 u}{\partial x^4} + \frac{\partial^2 u}{\partial x^2} + u\frac{\partial u}{\partial x} &= 0, \quad
    (x,t) \in (0,L) \times \bbR^+, \\[1ex]
    u(x,0) & = u_0, \quad \forall x \in [0,L].
    \end{align}
\end{subequations}
For the purposes of our experiments, the domain size is set to $L=32\pi$. Numerically, the KS equation \eqref{eq:KS_equation} is discretized in space using a Fourier pseudospectral method, which naturally enforces periodic boundary conditions. Time integration is carried out using the exponential time-differencing fourth-order Runge–Kutta (ETDRK4) scheme \cite{kassam2005fourth}. This method treats the stiff linear operator analytically, thereby avoiding the explicit linear Courant–Friedrichs–Lewy (CFL) constraint associated with fully explicit Runge–Kutta schemes. As a result, the time step is selected according to accuracy requirements and nonlinear resolution rather than linear stability limitations. We simulate the dynamics using a burn-in time of $200$ to ensure the trajectories are independent and sampled from the attractor. The observed trajectories are taken to be filtered solutions of \eqref{eq:KS_equation}, obtained by zeroing out the high Fourier modes of $u$. We denote this filtered solution by $\widetilde{u}$. In the smoothing setting we experiment with, the goal is to recover the full solution 
$u$ from $\widetilde{u}$, where $\widetilde{u}$ is obtained by retaining the first $64$ modes of $u$ and zeroing out the rest. In the forecasting setting, $\widetilde{u}$ is obtained by retaining the first $32$ modes of $u$ and zeroing out the rest.

For the smoothing experiment we use the observed trajectories on a time interval of $[0,100]$ as input data and predict the unobserved trajectories over the same time interval. We implement using points on the trajectories sampled at uniform intervals $\Delta t= 0.25$. For the forecasting experiment we use the observed trajectories on a time interval of $[0,100]$ as input data and predict the the trajectories over the future time $[100,102]$.
 We implement using points on the input and output trajectories sampled at uniform intervals $\Delta t= 0.25$. 
We train the neural operator solving the smoothing problem using $10000$ training trajectories and the one solving the forecasting problem using $80000$ trajectories. For both experiments we use a test set of $2000$ trajectories. We note the higher data complexity required to solve the forecasting problem. We summarize the parameter and experimental settings for the Kuramoto-Sivashinsky equation in \Cref{tab:ks_settings}.
\begin{table}[h]
\footnotesize
\centering
\caption{Kuramoto--Sivashinsky (KS) system settings. We recall that $u$ denotes the solution, while $\widetilde{u}$ is obtained by retaining the first $64$ modes of $u$ and zeroing out the rest. }
\label{tab:ks_settings}
\begin{tabular}{lcc}
    \toprule
    \textbf{Category} & \textbf{Smoothing} & \textbf{Forecasting} \\
    \midrule
    Parameters
        &\multicolumn{2}{c}{ $L = 32\pi$}
        \\
    \midrule
     Spatial Grid
        &\multicolumn{2}{c}{
          $x_j = jL/128,\quad j=1,\ldots,128$ %\quad
          %$ \bigl(u(x_1),\ldots,u(x_{128})\bigr) \in C([0,100];\bbR^{128})$
          }
          \\
    \midrule
    Observed
        & $ \bigl(\widetilde{u}(x_1),\ldots,\widetilde{u}(x_{128})\bigr) \!\in C([0,100];\bbR^{128})$
        & $ \bigl(\widetilde{u}(x_1),\ldots,\widetilde{u}(x_{128})\bigr)  \!\in C([0,100];\bbR^{128})$ \\
    \midrule
    Unobserved
        & $ \bigl(u(x_1),\ldots,u(x_{128})\bigr) \!\in C([0,100];\bbR^{128})$
        & $ \bigl(\widetilde{u}(x_1),\ldots,\widetilde{u}(x_{128})\bigr)  \!\in C([100,102];\bbR^{128})$ \\
    \midrule
    Obs. time
        &\multicolumn{2}{c}{ $\Delta t = 0.25$ }
        \\
    \bottomrule
\end{tabular}
\end{table}
In \Cref{fig:kse_smoothing} we demonstrate the qualitative results of the smoothing experiment. Indeed, we display the predicted trajectories corresponding to the median and highest relative $L^2$ errors in the test set compared to true trajectories. The panels showcase predictions that are qualitatively indistinguishable from the truth, indicating the success of the purely data-driven smoothing approach.
\begin{figure}[htbp]
    \centering
    \includegraphics[width=0.5\linewidth]{figures/KSE_median_worst_comparison.png}
    \caption{Median and worst-case performance on the KS system for smoothing.}
    \label{fig:kse_smoothing}
\end{figure}
In \Cref{fig:kse_spatiotemporal} we display results for the forecasting experiment. Here we show the pointwise errors relative to the root mean square of the ground truth (computed in space) for the sample in the test set achieving the median relative $L^2$ error. With the goal of obtaining forecasts on longer time horizons, we experiment with composition of the learned forecasting map $\Psi_F:C([0,100])\to C([100,102])$ to obtain $\Psi_F^n:C([0,100])\to C([100,100+2n])$. In \Cref{fig:kse_forecasting}, we show an example of predicted trajectories obtained by composing the map with $n=100$; the plot shows trajectory divergence after a handful of compositions, due to chaos and error compounding. However, in \Cref{fig:kse_forecasting_hist}, for $n=1000$ we compare the distribution of points in the predicted future trajectories $\widehat{p}$ to
the distribution of points in the ground truth trajectories $p$. The matching of distributions, and hence statistics, indicates the success in predicting valid trajectories from the attractor of the KS system, despite the chaotic nature of the equations and hence per-trajectory divergence. As further validation of the qualitative success of the purely data-driven forecasting approach, we recall the numerical comparison with the prediction given by the constant value of $p(T)$ for $T=100$, that is taking the predicted trajectory to be $\widehat{p}(t)=p(100)$ for $t\in[100,102]$. \Cref{tab:forecasting_results} shows the improvement in accuracy yielded by the transformer neural operator, in comparison to this naive approach.
\begin{figure}[htbp]
    \centering
    \includegraphics[width=0.6\linewidth]{figures/KSE_spatiotemporal_error.png}
    \caption{KS. Pointwise errors relative to the root mean square of the ground truth (computed in space) for the sample in the test set achieving the median relative $L^2$ error.}
    \label{fig:kse_spatiotemporal}
\end{figure}

\begin{figure}[htbp]
    \centering
    \begin{subfigure}{0.48\linewidth}
        \centering
        \includegraphics[width=\linewidth]{figures/KSE_median_worst_comparison_forecasting.png}
        \caption{}
        \label{fig:kse_forecasting}
    \end{subfigure}\hfill
    \begin{subfigure}{0.48\linewidth}
        \centering
        \includegraphics[width=\linewidth]{figures/KSE_histogram_comparison.png}
        \caption{}
        \label{fig:kse_forecasting_hist}
    \end{subfigure}
    \caption{KS. Left: forecast obtained by composing the learned map on one-step median and worst-case relative $L^2$ error samples. Right: distribution of trajectory predictions under composition of the learned forecasting map; the dark shade of red represents overlap between the blue ground truth distribution and the light red predicted distribution. The prevalence of the dark shade of red thus indicates matching of distributions.}
    \label{fig:kse_comparison}
\end{figure}

\section{Conclusions and Future Directions}
\label{sec:conclusions}

In this work we have formulated and analyzed a model for fully data-driven data assimilation, and
the smoothing and forecasting problems in particular; we have also defined neural operator algorithms to tackle the problem in practice. We show that in the context of dynamical systems, under an observability-rank condition, there exists a continuous operator mapping observations to unobserved or predicted quantities. We formulate universal approximation theorems  establishing the existence of neural operator parametrizations approximating the smoothing and forecasting operators up to arbitrary accuracy. We have further demonstrated these capabilities in practice by deploying transformer neural operators for these problems in the context of the Lorenz `63, Lorenz `96, and Kuramoto-Sivashinsky dynamical systems. 

The theory developed in this paper represents a step towards the analysis of data assimilation and forecasting methods that are model agnostic, avoiding the need for specified dynamics and possibly expensive model evaluations. This provides, for example, a first mathematical underpinning for the development of the direct-observation state estimators and direct-observation forecasts, which constitute the frontier of the AI-weather forecasting domain; see \cite{Allen2025} and \cite{alexe2024graphdop}, respectively. Nonetheless, many open challenges remain, 
and we highlight several interesting avenues for future work.
\begin{enumerate}
    \vspace{1ex}
    \item In \Cref{ex:Lorenz} we show that the observability-rank condition is satisfied by the Lorenz `63 dynamical system whenever $x_0\neq 0$. Characterizing for which points the Lorenz `96 dynamical system and Kuramoto-Sivashinsky equation satisfy this condition constitutes an avenue for further work. A similar analysis for the Navier-Stokes equation could lead to theory directly applicable to computational fluid dynamics and the atmospheric sciences.
    \vspace{1ex}
    \item In this paper we have not exploited any specific structural properties
    of the dynamics; but doing so might lead to more detailed understanding
    of the maps defining the smoothing and forecasting problems. For example
    the property of synchronization \cite{pecora1990synchronization} might be
    exploited, or the existence of an inertial manifold \cite{temam2012infinite}.
    %Several chaotic dynamical systems of interest in the data assimilation context exhibit the property of synchronization \cite{pecora1990synchronization}. Synchronization of trajectories of coupled chaotic dynamical systems occurs, for example, in the setting of \eqref{eq:dynamical_system_intro} supplemented with additive linear terms $A_pp$ and $A_qq$ for negative definite $A_q$, and with Lipschitz conditions on the functions $f,g$. An investigation of the interplay between the universal approximation properties of neural operators and the synchronization properties of the dynamics constitutes an interesting avenue for future work.
    \vspace{1ex}
    \item As mentioned in \Cref{rem:global} it is of interest to extend the local existence results in this paper to show under a different set of assumptions the existence of an operator mapping observations to unobserved variables, defined on the whole domain of the dynamics. Applying some of the ideas in \cite{sauer1991embedology} constitutes a promising avenue for future work.
    \vspace{1ex}
    \item It is of interest to perform a detailed numerical study comparing the performance of different neural operator architectures when applied to the data assimilation problems presented in this paper. Furthermore, an accuracy-complexity tradeoff analysis between neural operators and non learning-based smoothing algorithms would be of interest.
    \vspace{1ex}
    \item It is of interest to study discrete time dynamical systems from the perspective developed in this paper, and to use this setting to make concrete connections to Takens' embedding theorem \cite{takens2006detecting,sauer1991embedology}.
\end{enumerate}

%\mm{The main result we have shown gives us a local condition that allows to recover initial conditions from partial observation. However, it is not possible to patch these local conditions into a global operator since we can't ``glue the inverse maps together'' (I can explain if needed, but this is a standard result and I think can/should be omitted). However, we can get a more global result to allow us to learn long term dynamics of the attractor as follows. From Fractal Delay Embedding Theorem Thm. 2.5 from Embedology Sauer, Yorke, Casdagli, we can consider the partially observed data to be a continuous version of a delay embedding map (i.e. if we sampled at equally spaced time intervals from some time spacing $\tau$ we would get a delay-embedding map). For almost every (in a prevalent set view, so this means if ``perturb the map with respect to some dense set" it holds for some map in the ball) delay embedding map, we have that these partial observations distinguish points (i.e. initial conditions) sampled from the attractor. This is relevant to us because if we take a compact set of the attractor we can learn the inverse map from the observed data to the initial condition. The coordinate delay map is a smooth map (or sufficient regularity, at least continuous). If you have a continuous map over an injective domain then it is continuously invertible. I.e. we can find a continuous map that takes us from the observed partial data to the initial condition from the attractor. integrating forward gets us the smoothing map and the forecasting map. This is the generalized W for systems with an attractor and gives us justification for learning big sets not just neighborhoods of points.}



\appendix

\section{Observability-Rank Condition in Control Theory} 
\label{appendix:observability_control}

In the following discussion we explain the connection between the observability-rank condition, employed in this paper, and the observability-rank condition from control theory, introduced in the seminal paper \cite{hermann1977nonlinear}. In the general control-theoretic context, consider the partially observed dynamical system
\begin{align}
\label{eq:hermann_ds}
    \begin{aligned}
        \dot{x} &= \phi(x, u),\\
        y &= h(x),
    \end{aligned}
\end{align}
where $u(t) \in \Omega\subset\R^l$ for each $t \in \R^+$, $x \in M$, a $C^{\infty}$ connected manifold of dimension $m$, $y \in \R^n$, and $\phi,h$ are $C^{\infty}$ functions. For simplicity, in what follows we
consider the setting where $M=\R^{d_p + d_q}$ and $h: \R^{d_p+d_q} \to \R^{d_p}$. The following links the main body of the paper to this setting.


\begin{example} \label{ex:ex1}
The partially observed dynamical system \eqref{eq:dynamical_system_intro} we consider in this paper may be written in the form
\begin{align}
\label{eq:hermann_ds_ours}
    \begin{aligned}
        \dot{p} &= f(p,q),\\
        \dot{q} &= g(p,q),\\
        p &= \pi_p(p,q),
    \end{aligned}
\end{align}
This is a specific example of the control theoretic setting with $M=\R^{d_p + d_q}$, the external control $u$ absent, $x=(p,q)$, $\phi=(f,g)$ and $h = \pi_p$.
\end{example}

We now introduce the technical notions required to define the control theoretic
version of the observability-rank condition from \cite{hermann1977nonlinear}; these definitions may be found 
in \cite{hermann1977nonlinear} and are developed for general $C^\infty$ manifolds $M$ and general $C^\infty$ observation operators $h:\R^m\to \R^n$. Here we work
in the setting of $M=\R^{d_p + d_q}$ and $h:\R^{d_p+d_q}\to \R^{d_p}$, aligning with the Example \ref{ex:ex1},
and in particular with the contents of our paper.

Every point $\mathsf{u} \in \Omega$ defines a vector field $x\mapsto \phi(x, \mathsf{u}) \in C^\infty(\R^{d_p+d_q};\R^{d_p+d_q})$; we let $\mathcal{F}^{0}\subset C^\infty(\R^{d_p+d_q};\R^{d_p+d_q})$ denote the collection of all of these vector fields, one for each $\mathsf{u} \in \Omega$.
We let $\mathcal{H}^{0}$ denote the subset of $C^\infty(\R^{d_p+d_q})$ consisting of the $d_p$ component functions $h_1, h_2, \ldots, h_{d_p}:\R^{d_p+d_q} \to \R$ of the observation map $h:\R^{d_p+d_q}\to \R^{d_p}$, viewed as scalar-valued smooth functions on $\R^{d_p+d_q}$. We define $\mathcal{H}$ to be the smallest linear subspace of $C^\infty(\R^{d_p+d_q})$ containing $\mathcal{H}^{0}$ and which is closed with respect to Lie differentiation {under} elements of $\mathcal{F}^{0}$. For each $\varphi \in \mathcal{H}$, the differential $\sD\varphi(x) : \mathbb{R}^{d_p+d_q} \to \mathbb{R}$ is a linear functional on $\mathbb{R}^{d_p+d_q}$, i.e.\ an element of $(\mathbb{R}^{d_p+d_q})^*$. We define $\mathrm{d}\mathcal{H}^0 = \{\sD\varphi : \varphi \in \mathcal{H}^0\}$ and $\mathrm{d}\mathcal{H} = \{\sD\varphi : \varphi \in \mathcal{H}\}$ as collections of maps $\mathbb{R}^{d_p+d_q} \to (\mathbb{R}^{d_p+d_q})^*$, and $\mathrm{d}\mathcal{H}(x) \subseteq (\mathbb{R}^{d_p+d_q})^*$ as the subspace obtained by evaluating elements of $\mathrm{d}\mathcal{H}$ at $x$.



\begin{definition}[Observability-rank condition in \cite{hermann1977nonlinear}]
    The dynamical system \eqref{eq:hermann_ds} satisfies the \textit{observability-rank condition} at $\fx$ if $\dim(\rd\mathcal{H}(\fx)) = d_p+d_q$.
\end{definition}


With the translation between our paper and \cite{hermann1977nonlinear}, as described in Example \ref{ex:ex1}, we may formulate the following proposition.

\begin{proposition}
\label{prop:connection_control}
    The dynamical system \eqref{eq:dynamical_system_intro} satisfying \Cref{assump:regularity} with $f,g\in C^\infty$ and satisfying \Cref{assump:observability} at the point $(\fp,\fq)$ is an instance of a dynamical system \eqref{eq:hermann_ds} with a fixed control, satisfying the observability-rank condition from \cite{hermann1977nonlinear} at the point $\fx=(\fp,\fq)$.
\end{proposition}



\begin{proof}[Proof of \Cref{prop:connection_control}]
We are in the specific setting of Example \ref{ex:ex1} and hence, because no control is present,
the set $\mathcal{F}^0$ simply comprises the vector field $(f,g)\in C^\infty(\R^{d_p+d_q};\R^{d_p+d_q})$. 
Letting $\fx = (\fp, \fq)$ and recalling the notation for Lie derivatives defined in \Cref{subsec:notation}, we have that
\begin{subequations}
\label{eq:LiePi}
\begin{align}
    \mathcal{L}_{(f,g)}(\pi_p)(\fx) &= \frac{\partial \pi_p}{\partial x}(\fx)\begin{pmatrix}
        f(\fx)\\
        g(\fx)
    \end{pmatrix}\\
    &= \begin{pmatrix}
        I_{d_p \times d_p} & \mathbf{0}_{d_p \times d_q}
    \end{pmatrix}\begin{pmatrix}
        f(\fp, \fq)\\
        g(\fp, \fq)
    \end{pmatrix} = f(\fp, \fq).
\end{align}
\end{subequations}
Therefore, higher order Lie derivatives of $\pi_p$ under the vector field $(f,g)$ yield $\cL^i_{(f,g)}\, f$. Let $\mathcal{G}$ be the span of the component functions 
$\{\pi_{p,1}, \ldots, \pi_{p,d_p}\} \cup \{\mathcal{L}^i_{(f,g)} f_j : 1 \leq j \leq d_p,\, 0 \leq i \leq n-1\}$,
where $\pi_{p,j} : \mathbb{R}^{d_p+d_q} \to \mathbb{R}$ denotes the $j$th component of $\pi_p$ and similarly $f_j$ denotes the $j$th component of $f$, so that the evaluations $\{\pi_{p,j}(\fp,\fq)\}_{j=1}^{d_p}$ and $\{\mathcal{L}^i_{(f,g)} f_j(\fp,\fq)\}_{j=1}^{d_p}$ correspond to the $(n+1)d_p$ entries of $F^{(n)}(\fp,\fq)$ in \eqref{eq:Lie}.

From \eqref{eq:LiePi} and closure of $\mathcal{H}$ under Lie differentiation under $(f,g)$, we have that $\cG \subseteq \mathcal{H}$; consequently, $\rd\mathcal{H}$ contains differentials of elements of $\cG$. The linear transformation $L:\R^{(n+1)d_p}\to \R^{d_p+d_q}$ from \eqref{eq:observability_eq_inv2} defines a map 
\begin{align}
    LF^{(n)} \colon (\fp, \fq) \mapsto [s_1(\fp, \fq), s_2(\fp,\fq), \ldots, s_{d_p + d_q}(\fp, \fq)]^{\top},
\end{align}
for $s_i \in \cG$. Since $\mathrm{d}\mathcal{H}(\fp,\fq) \subseteq (\mathbb{R}^{d_p+d_q})^*$, we have that $\dim(\rd\mathcal{H}(\fp, \fq)) \leq d_p+d_q$. Since $s_1, \ldots, s_{d_p+d_q} \in \mathcal{G} \subseteq \mathcal{H}$, their differentials 
$\sD s_1(\fp,\fq), \ldots, \sD s_{d_p+d_q}(\fp,\fq)$ are elements of $\mathrm{d}\mathcal{H}(\fp,\fq)$. 
Moreover, these differentials are the rows of the Jacobian $\sD(LF^{(n)})(\fp,\fq)$, 
which has rank $d_p+d_q$ by \Cref{assump:observability}, and hence are linearly independent. 
Therefore $\dim(\mathrm{d}\mathcal{H}(\fp,\fq)) \geq d_p+d_q$, and combined with the upper bound 
$\dim(\mathrm{d}\mathcal{H}(\fp,\fq)) \leq d_p+d_q$, we conclude that $\dim(\mathrm{d}\mathcal{H}(\fp,\fq)) = d_p+d_q$.
\end{proof}

\section{Architectural Details}
\label{appendix:architectures}

In this section we describe the transformer neural operator architectures employed in \Cref{sec:experiments},
highlighting the use of self- and cross-attention. The goal of this section is two-fold: to present the architectural details, but also to discuss how the architecture may be cast in the universal approximation context of \Cref{thm:UANeuralOperator}; recall that this proposition is used in the proof of \Cref{thm:UAsmoothing,thm:UAforecasting} which underpin the methodology developed and deployed in this paper. The first transformer neural operator uses self-attention. The second uses a combination of self-attention and cross-attention. The use of cross-attention enables the output of the architecture to be defined on a different number of grid points than the input. The self-attention based transformer neural operator we use is described in detail in \cite[Section 4.2]{calvello2024continuum}. A slight modification to the architecture implemented in practice leads to a universal approximation theorem of the form we apply in this paper (as in \Cref{thm:UANeuralOperator}), the statement of which may be found in \cite[Theorem 22]{calvello2024continuum}. We focus the following discussion on the cross-attention based transformer neural operator. To this end, we recall the definition of cross-attention in function space from \cite{calvello2024continuum}.

Let $D \subseteq \R^d$ and $E \subseteq \R^e$ be open sets. Let $u: D \to \R^{d_u}$ be a function,
$v: E \to \R^{d_v}$ another function, and $x \in E$, $y \in D$ points. Then, for learnable parameters $Q\in \R^{d_K\times d_v},K\in \R^{d_K\times d_u},V\in \R^{d_V\times d_u}$, cross-attention is defined as the operator acting on functions $v,u$ such that
\begin{equation}
\label{eq:cross_attention}
    \sC(v,u)(x) = \bbE_{y \sim \pi(y;v,u,x)} [V u(y)],
\end{equation}
for any $x\in E$, where the probability density function $\pi(\cdot\,;v,u,x): D \to \R^+$ is defined as
\begin{equation}
\label{eq:cross_attention_pdf}
    \pi(y; v,u, x) = \frac{\mathrm{exp} \Big ( \big\langle Q v(x), K u(y) \big\rangle_{\R^{\dK}} \Big )}{\int_{D} \mathrm{exp} \Big ( \big\langle Q v(x), K u(s) \big\rangle_{\R^{\dK}} \Big ) \: \mathrm{d}s},
\end{equation}
for any $y \in {D}$. We note that self-attention is a special case of cross-attention where
$D=E$ and $v=u$. In the following we let $\cV,\cU,\cW,\cZ$ denote arbitrary function spaces; furthermore, while cross-attention allows a more general definition, here we write the architecture as a mapping between functions over the common domain $D$. The cross-attention based transformer neural operator we employ may be written as the map $\Psi_\sC(\placeholder, \placeholder;\theta):\cV(D;\R^\ell)\times \cU(D;\R^r)\to \cZ(D;\R^{r'})$ acting on functions $v\in \cV(D;\R^{\ell})$ and $u\in \cU(D;\R^r)$ such that
\begin{equation}
\label{eq:transformer_operator_compact}
\Psi_{\sC}(v, u;\theta) \coloneqq \Bigl(\sT_{\textrm{out}}\circ\sD(v, \placeholder)\circ\sE_L \circ \sT_{\textrm{in}}\Bigr) (u,\theta).
\end{equation}
In practice, we define $\sT_{\textrm{in}}:\cU(D;\R^r)\to \cW(D;\R^c)$ via concatenation with the grid (positional encoding) and application of a learnable linear transformation that lifts to a latent channel dimension $c$. The operator $\sE_L: \cW(D;\R^c)\to \cW(D;\R^c)$ is defined via a composition of $L$ self-attention layers. Its construction is outlined in detail in \cite[Section 4]{calvello2024continuum}. The transformer decoder $\sD: \cV(D;\R^\ell)\times \cW(D;\R^c)\to \cW(D;\R^c)$ is defined for $(v,u)\in \cV(D;\R^\ell)\times \cW(D;\R^c)$ via the iteration
\begin{subequations}
\label{eq:recurrence_enc_def}
\begin{align}
    v &\mapsfrom S_{\mathrm{in}}v,\\
    v &\mapsfrom {W^\sD_1}v + \sC(v, v), \\
    v &\mapsfrom \sF_{\textrm{LayerNorm}}(v),\\
    \label{eq:reason}
    v &\mapsfrom {W^\sD_2}v + \sC(v, u), \\
    v &\mapsfrom \sF_{\textrm{LayerNorm}}(v), \\
    v &\mapsfrom {W^\sD_3}v + \sF_{\textrm{NN}} (v),\\
    v &\mapsfrom \sF_{\textrm{LayerNorm}}(v),
\end{align}
\end{subequations}
where $S_{\mathrm{in}}$ is a learnable linear transformation lifting to the channel dimension, $\sF_{\textrm{LayerNorm}}$ denotes application of layer normalization, $\sF_{\textrm{NN}}$ a two layer multilayer perceptron. We note that in practice, multihead cross-attention is employed; we refer to \cite[Section 4]{calvello2024continuum} for a discussion on multihead attention. We also note that in our implementation, we fix the linear transformation $W^{\sD}_1, W^{\sD}_2, W^{\sD}_3$ to be the identity. For implementation purposes, we define the operator $\sT_{\textrm{out}}:\cW(D;\R^c)\to \cZ(D;\R^{r'})$ as a linear transformation projecting to the output space.

Due to the definition of cross-attention and step \eqref{eq:reason} in the decoder, the discretization size of the output function corresponds to the discretization size of the input function $v$, making this architecture particularly attractive for forecasting problems, where the output time interval may be of length differing to the input. In practical application of this architecture, the $v$ is fixed to be the rescaled identity function, sampled at the grid points where the output function is to be evaluated. Since the $v$ is fixed, for the approximation theory developed in the following section we only require consideration of the operator $\Psi(v,\placeholder;\theta): \cU(D;\R^r)\to \cZ(D;\R^{r'})$. Because the architecture is the discrete version of the continuum neural operator, the scheme is discretization invariant in both inputs, meaning that the parameters will be independent of any discretization of input and output functions; the resulting scheme is thus deployable at any discretization. In the next subsection we state and prove a universal approximation theorem for a slight variant of the scheme implemented in practice.

\subsection{Universal Approximation for Cross-Attention}

Consider activation functions $\sigma \in C^\infty(\R)$ which are non-polynomial and Lipschitz continuous. We consider neural operators of the form \eqref{eq:transformer_operator_compact} where $\sT_{\textrm{in}}: \cU(D;\R^r)\to \cW(D;\R^c)$ and $\sT_{\textrm{out}}: \cW\bigl(D;\R^c\bigr)\to\cZ\bigl(D;\R^{d_z} \bigr)$ are defined by neural networks of the form 
\begin{align}
    \Bigl(\sT_{\textrm{in}}(u)\Bigr)(x) &= R_2\sigma\Bigl(R_1\bigl(u(x),x\bigr)+b_R \Bigr)+b_R',\\
    \Bigl(\sT_{\textrm{out}}(v)\Bigr)(x) &= P_2\sigma\Bigl(P_1\bigl(v(x),x\bigr)+b_P \Bigr)+b_P',
\end{align}
where $\bigl(u(x),x\bigr)\in \R^{r+d}$ and $\bigl(v(x),x\bigr)\in \R^{c+d}$, where $R_1,R_2,P_1,P_2$ are learned linear transformations of appropriate dimensions and $b_R,b_R',b_P,b_P'$ are learned vectors. We define the operator $\sE: \cW\bigl(D;\R^c \bigr) \to \cW\bigl(D;\R^c \bigr)$ as a two-step map acting on its inputs $u\in\cW(D;\R^c)$ as
\begin{subequations}
\label{eq:recurrence_enc_func_ua}
\begin{align}
    {u}(x) &\mapsfrom W_1^\sE{u}(x) + \sC\big({u},u; Q^\sE,K^\sE,V^\sE\big)(x), \\
    u(x) &\mapsfrom W_2^\sE{u}(x) + W_3^\sE \sigma\bigl(W_4^\sE{u}(x)+b_1^\sE\bigr) + b_2^\sE,
\end{align}
\end{subequations}
for any $x\in D$. Note that the layer thus defined does not include layer normalizations, and hence is a variant of the self-attention transformer layer employed in practice. We also define a no-layer normalization variant of the decoder in \eqref{eq:recurrence_enc_def}, which also includes a residual of the output of the encoder after application of cross-attention; indeed we consider $\sD: \cV(D;\R^\ell)\times \cW(D;\R^c)\to \cW(D;\R^c)$ defined for $(v,u)\in \cV(D;\R^\ell)\times \cW(D;\R^c)$ via the iteration
\begin{subequations}
\label{eq:recurrence_enc_def1}
\begin{align}
    v(x) &\mapsfrom S_{\mathrm{in}}v(x)\\
    v(x) &\mapsfrom {W^\sD_1}v(x) + \sC(v, v; Q^\sD_1,K^\sD_1,V^\sD_1)(x), \\
    v(x) &\mapsfrom {W^\sD_2}\bigl(v(x),u(x)\bigr) + \sC(v, u; Q^\sD_2,K^\sD_2,V^\sD_2)(x), \\
    v(x) &\mapsfrom {W^\sD_3}v(x) + W_4^\sD \sigma\bigl(W_5^\sD{v}(x)+b_1^\sD\bigr)+ b_2^\sD,
\end{align}
\end{subequations}
for any $x\in D$. We may now apply the result of \cite{lanthaler2023nonlocal} to show two universal approximation theorems for the resulting cross-attention based transformer neural operator.


\begin{theorem}
\label{thm:ua1}
Let $D\subset \R^d$ be a bounded domain with Lipschitz boundary, and fix integers $s,s'\geq 0$. If $\Psi^\dagger:C^s\bigl(\Bar{D};\R^r\bigr) \to C^{s'}\bigl(\Bar{D};\R^{r'}\bigr)$
is a continuous operator and $K\subset C^s\bigl(\Bar{D};\R^r\bigr)$ a compact set, then for any $\epsilon>0$, there exists a cross-attention based transformer neural operator 
$\Psi(v,\placeholder;\theta):K\subset C^s\bigl(\Bar{D};\R^r\bigr) \to C^{s'}\bigl(\Bar{D};\R^{r'}\bigr)$ so that
\begin{equation}
\label{eq:ua_condition}
\sup_{u\in K} \left\|\Psi^\dagger(u) - \Psi(v,u;\theta)   \right\|_{C^{s'}} \leq \epsilon.
\end{equation}
\end{theorem}

\begin{proof}
We begin by noting that for $Q^\sD_2,K^\sD_2=0$ and $V^\sD_2 = I$, the cross-attention mapping employed in the decoder reduces to 
\begin{equation}
\sC\bigl(v,u;Q^\sD_2,K^\sD_2,V^\sD_2 \bigr)(\placeholder) = \mathbb{E}_{y\sim \pi(y;v,u,\placeholder)}[V^\sD_2 u(y)] = \frac{1}{|D|}\int u(x)~\mathrm{d} x.
\end{equation}
For encoder weights $V^\sD=0$, $W_3^\sE=0$, $W_1^\sE=W_2^\sE=I$ and decoder weights $W^\sD_2=\bigl(0, W\bigr)$, $W^\sD_3 = 0$, $W^\sD_4=W^\sD_5=I$ and $b^\sD_2 = 0$, the composition of encoder \eqref{eq:recurrence_enc_func_ua} and decoder \eqref{eq:recurrence_enc_def1} reduces to the mapping 
\begin{equation}
\label{eq:nonlocal_NO}
    u(\placeholder) \mapsto \sigma\left(Wu(\placeholder) + b_1 + \frac{1}{|D|}\int u(x)~\mathrm{d} x \right).
\end{equation}
The existence of $W,R_1,R_2,P_1,P_2,b_1,b_R,b_R',b_P,b_P'$ so that $\Psi(v,\placeholder\,;\theta)$ satisfies \eqref{eq:ua_condition} then follows from \cite[Theorem 2.1]{lanthaler2023nonlocal}, which also involves the application of the universality result for two-layer neural networks of \cite[Theorem 4.1]{Pinkus_1999}.
\end{proof}


\section*{Acknowledgments}

AMS is supported by a Department of Defense (DoD) Vannevar Bush Faculty Fellowship (award N00014-22-1-2790).
Both EC's acknowledges support from the Vannevar Bush Faculty Fellowship held by AS; Calvello also
acknowledges support from the Resnick Sustainability Institute and Carlson the Caltech Von Karman Instructorship.


\bibliographystyle{siamplain}
\bibliography{references}

\end{document}

